% !TeX program = xelatex
\documentclass[11pt,a4paper]{article}

% ============================================================================
% atl-keys.tex -- configuration keys
%
%   \atlset{key=value, ...}     set (or override) keys; any key name is allowed
%   \atlmeta{key}               value of a key (empty if unset)
%   \atlifset{key}{yes}{no}     is the key set and non-empty?
%   \atlifeq{key}{val}{y}{n}    does the key equal val?
%   \atlifon{key}{yes}{no}      is the key true/on/yes?
%   \atldefault{key}{value}     set the key only if it is not set yet
%   \atlstr{name}               localised label (defined by the language layer)
%
% Order in main.tex: this file, template defaults, config.tex (your values), atl.tex.
% ============================================================================
\makeatletter
\RequirePackage{pgfkeys}
\RequirePackage{etoolbox}
\RequirePackage{pdftexcmds}

\pgfkeys{/atl/.unknown/.code={\expandafter\gdef\csname atl@m@\pgfkeyscurrentname\endcsname{#1}}}

\newcommand{\atlset}[1]{\pgfkeys{/atl/.cd,#1}}
\newcommand{\atlmeta}[1]{\ifcsname atl@m@#1\endcsname\csname atl@m@#1\endcsname\fi}

% The conditionals select a branch first and only then run it, so the chosen code may
% contain & and \\ (table cells) without sitting inside an open \if.
\long\def\atl@iftrue#1#2{#1}\long\def\atl@iffalse#1#2{#2}
\newcommand{\atl@isset}[1]{%
  \ifcsname atl@m@#1\endcsname
    \expandafter\ifx\csname atl@m@#1\endcsname\@empty false\else true\fi
  \else false\fi}
\newcommand{\atl@iseq}[2]{%
  \ifcsname atl@m@#1\endcsname
    \ifnum\pdf@strcmp{\csname atl@m@#1\endcsname}{#2}=\z@ true\else false\fi
  \else false\fi}
\newcommand{\atlifset}[3]{\csname atl@if\atl@isset{#1}\endcsname{#2}{#3}}
\newcommand{\atlifeq}[4]{\csname atl@if\atl@iseq{#1}{#2}\endcsname{#3}{#4}}
\newcommand{\atlifon}[3]{%
  \atlifeq{#1}{true}{#2}{\atlifeq{#1}{on}{#2}{\atlifeq{#1}{yes}{#2}{#3}}}}
\newcommand{\atldefault}[2]{\atlifset{#1}{}{\atlset{#1={#2}}}}

% Localised strings: defined with \atldefstr in atl-lang-*.tex
\newcommand{\atldefstr}[2]{\expandafter\gdef\csname atl@s@#1\endcsname{#2}}
\newcommand{\atlstr}[1]{\ifcsname atl@s@#1\endcsname\csname atl@s@#1\endcsname\else[#1]\fi}
\makeatother
                      % ۱. کلیدهای پیکربندی
\atlset{language=fa, category=homework,   % ۲. مشخصات این قالب (آزادانه تغییر دهید)
        titlepage=block, margins=normal, linespread=1.2,
        infofields={author,studentid,course,instructor,duedate},
        title={تمرین}, date={}}
% Written by mktex on 2026-10-17. Edit freely; anything set here overrides the template defaults.
\atlset{
  language = fa,
  palette  = blue,
  listing  = darcula,
  title = {تمرین سری اول – بخش تئوری},
  course = {سیستم‌عامل},
  instructor = {دکتر اسدی},
  term = {نیم‌سال اول ۰۶-۰۵},
  author = {محمدحسین سبزعلیان},
  studentid = {402101906},
}
                            % ۳. مقادیر شما (به‌وسیلهٔ mktex نوشته می‌شود)
% ============================================================================
% atl.tex -- AdvancedLatexTemplate engine loader
%
% main.tex of every template:
%     \documentclass[...]{article}            (or report)
%     % ============================================================================
% atl-keys.tex -- configuration keys
%
%   \atlset{key=value, ...}     set (or override) keys; any key name is allowed
%   \atlmeta{key}               value of a key (empty if unset)
%   \atlifset{key}{yes}{no}     is the key set and non-empty?
%   \atlifeq{key}{val}{y}{n}    does the key equal val?
%   \atlifon{key}{yes}{no}      is the key true/on/yes?
%   \atldefault{key}{value}     set the key only if it is not set yet
%   \atlstr{name}               localised label (defined by the language layer)
%
% Order in main.tex: this file, template defaults, config.tex (your values), atl.tex.
% ============================================================================
\makeatletter
\RequirePackage{pgfkeys}
\RequirePackage{etoolbox}
\RequirePackage{pdftexcmds}

\pgfkeys{/atl/.unknown/.code={\expandafter\gdef\csname atl@m@\pgfkeyscurrentname\endcsname{#1}}}

\newcommand{\atlset}[1]{\pgfkeys{/atl/.cd,#1}}
\newcommand{\atlmeta}[1]{\ifcsname atl@m@#1\endcsname\csname atl@m@#1\endcsname\fi}

% The conditionals select a branch first and only then run it, so the chosen code may
% contain & and \\ (table cells) without sitting inside an open \if.
\long\def\atl@iftrue#1#2{#1}\long\def\atl@iffalse#1#2{#2}
\newcommand{\atl@isset}[1]{%
  \ifcsname atl@m@#1\endcsname
    \expandafter\ifx\csname atl@m@#1\endcsname\@empty false\else true\fi
  \else false\fi}
\newcommand{\atl@iseq}[2]{%
  \ifcsname atl@m@#1\endcsname
    \ifnum\pdf@strcmp{\csname atl@m@#1\endcsname}{#2}=\z@ true\else false\fi
  \else false\fi}
\newcommand{\atlifset}[3]{\csname atl@if\atl@isset{#1}\endcsname{#2}{#3}}
\newcommand{\atlifeq}[4]{\csname atl@if\atl@iseq{#1}{#2}\endcsname{#3}{#4}}
\newcommand{\atlifon}[3]{%
  \atlifeq{#1}{true}{#2}{\atlifeq{#1}{on}{#2}{\atlifeq{#1}{yes}{#2}{#3}}}}
\newcommand{\atldefault}[2]{\atlifset{#1}{}{\atlset{#1={#2}}}}

% Localised strings: defined with \atldefstr in atl-lang-*.tex
\newcommand{\atldefstr}[2]{\expandafter\gdef\csname atl@s@#1\endcsname{#2}}
\newcommand{\atlstr}[1]{\ifcsname atl@s@#1\endcsname\csname atl@s@#1\endcsname\else[#1]\fi}
\makeatother
                    keys
%     \atlset{category=..., palette=..., ...} this template's defaults
%     % Written by mktex on 2026-10-17. Edit freely; anything set here overrides the template defaults.
\atlset{
  language = fa,
  palette  = blue,
  listing  = darcula,
  title = {تمرین سری اول – بخش تئوری},
  course = {سیستم‌عامل},
  instructor = {دکتر اسدی},
  term = {نیم‌سال اول ۰۶-۰۵},
  author = {محمدحسین سبزعلیان},
  studentid = {402101906},
}
                          your values (written by mktex)
%     % ============================================================================
% atl.tex -- AdvancedLatexTemplate engine loader
%
% main.tex of every template:
%     \documentclass[...]{article}            (or report)
%     % ============================================================================
% atl-keys.tex -- configuration keys
%
%   \atlset{key=value, ...}     set (or override) keys; any key name is allowed
%   \atlmeta{key}               value of a key (empty if unset)
%   \atlifset{key}{yes}{no}     is the key set and non-empty?
%   \atlifeq{key}{val}{y}{n}    does the key equal val?
%   \atlifon{key}{yes}{no}      is the key true/on/yes?
%   \atldefault{key}{value}     set the key only if it is not set yet
%   \atlstr{name}               localised label (defined by the language layer)
%
% Order in main.tex: this file, template defaults, config.tex (your values), atl.tex.
% ============================================================================
\makeatletter
\RequirePackage{pgfkeys}
\RequirePackage{etoolbox}
\RequirePackage{pdftexcmds}

\pgfkeys{/atl/.unknown/.code={\expandafter\gdef\csname atl@m@\pgfkeyscurrentname\endcsname{#1}}}

\newcommand{\atlset}[1]{\pgfkeys{/atl/.cd,#1}}
\newcommand{\atlmeta}[1]{\ifcsname atl@m@#1\endcsname\csname atl@m@#1\endcsname\fi}

% The conditionals select a branch first and only then run it, so the chosen code may
% contain & and \\ (table cells) without sitting inside an open \if.
\long\def\atl@iftrue#1#2{#1}\long\def\atl@iffalse#1#2{#2}
\newcommand{\atl@isset}[1]{%
  \ifcsname atl@m@#1\endcsname
    \expandafter\ifx\csname atl@m@#1\endcsname\@empty false\else true\fi
  \else false\fi}
\newcommand{\atl@iseq}[2]{%
  \ifcsname atl@m@#1\endcsname
    \ifnum\pdf@strcmp{\csname atl@m@#1\endcsname}{#2}=\z@ true\else false\fi
  \else false\fi}
\newcommand{\atlifset}[3]{\csname atl@if\atl@isset{#1}\endcsname{#2}{#3}}
\newcommand{\atlifeq}[4]{\csname atl@if\atl@iseq{#1}{#2}\endcsname{#3}{#4}}
\newcommand{\atlifon}[3]{%
  \atlifeq{#1}{true}{#2}{\atlifeq{#1}{on}{#2}{\atlifeq{#1}{yes}{#2}{#3}}}}
\newcommand{\atldefault}[2]{\atlifset{#1}{}{\atlset{#1={#2}}}}

% Localised strings: defined with \atldefstr in atl-lang-*.tex
\newcommand{\atldefstr}[2]{\expandafter\gdef\csname atl@s@#1\endcsname{#2}}
\newcommand{\atlstr}[1]{\ifcsname atl@s@#1\endcsname\csname atl@s@#1\endcsname\else[#1]\fi}
\makeatother
                    keys
%     \atlset{category=..., palette=..., ...} this template's defaults
%     % Written by mktex on 2026-10-17. Edit freely; anything set here overrides the template defaults.
\atlset{
  language = fa,
  palette  = blue,
  listing  = darcula,
  title = {تمرین سری اول – بخش تئوری},
  course = {سیستم‌عامل},
  instructor = {دکتر اسدی},
  term = {نیم‌سال اول ۰۶-۰۵},
  author = {محمدحسین سبزعلیان},
  studentid = {402101906},
}
                          your values (written by mktex)
%     % ============================================================================
% atl.tex -- AdvancedLatexTemplate engine loader
%
% main.tex of every template:
%     \documentclass[...]{article}            (or report)
%     \input{atl/atl-keys}                    keys
%     \atlset{category=..., palette=..., ...} this template's defaults
%     \input{config}                          your values (written by mktex)
%     \input{atl/atl}                         <- this file (also loads preamble/packages.tex)
%     \input{preamble/local}                  your own macros and settings
%     \begin{document} \atlmaketitle ...
% Needs XeLaTeX:  latexmk -xelatex main.tex
% ============================================================================
\makeatletter
\atldefault{language}{en}
\atldefault{category}{homework}
\newif\ifatlfa
\atlifeq{language}{fa}{\atlfatrue}{\atlfafalse}

\input{atl/atl-packages}
\IfFileExists{preamble/packages.tex}{\input{preamble/packages.tex}}{}  % your extra \usepackage lines (before XePersian)
\ifatlfa\input{atl/atl-lang-fa}\else\input{atl/atl-lang-en}\fi
\input{atl/atl-palettes}
\input{atl/atl-layout}
\input{atl/atl-headings}
\input{atl/atl-pagestyle}
\input{atl/atl-floats}
\input{atl/atl-boxes}
\input{atl/atl-code}
\input{atl/atl-macros}
\input{atl/atl-theorems}
\input{atl/atl-titles}
\input{atl/atl-bib}

\makeatletter % (each module ends with \makeatother)
% names LaTeX prints itself (Contents, List of Figures, ...) come from the language layer
\AtBeginDocument{%
  \renewcommand{\contentsname}{\atlstr{contents}}%
  \renewcommand{\listfigurename}{\atlstr{listoffigures}}%
  \renewcommand{\listtablename}{\atlstr{listoftables}}%
  \renewcommand{\figurename}{\atlstr{figure}}%
  \renewcommand{\tablename}{\atlstr{table}}%
  \renewcommand{\appendixname}{\atlstr{appendix}}%
  \@ifundefined{refname}{\newcommand{\refname}{\atlstr{references}}}{\renewcommand{\refname}{\atlstr{references}}}%
  \@ifundefined{bibname}{}{\renewcommand{\bibname}{\atlstr{references}}}%
  \@ifundefined{abstractname}{}{\renewcommand{\abstractname}{\atlstr{abstract}}}}
\hypersetup{linktoc=page}
\makeatother
\makeatother
                         <- this file (also loads preamble/packages.tex)
%     % ماکروها و تنظیمات شخصی شما؛ پس از موتور قالب بارگذاری می‌شود.
% \atlsetpalette{green}                 % تغییر پالت رنگ از همین‌جا
% \atlsetlisting{bright}                % تغییر پوستهٔ کد (darcula, dark, bright, github)
% \newcommand{\A}{\mat{A}}

% ---- ماکروهای این تمرین ----
\definecolor{navy}{HTML}{1F3A68}
\definecolor{accent}{HTML}{B3401E}
\definecolor{lane1}{HTML}{E8F0FE}
\definecolor{lane2}{HTML}{FDF1E3}
\newcommand{\code}[1]{\lr{\texttt{#1}}}                      % نام فراخوانی، ثبات، مسیر
\newcommand{\src}[1]{\par\begingroup\color{gray}\small\noindent\textbf{منبع:}~#1\par\endgroup}
                  your own macros and settings
%     \begin{document} \atlmaketitle ...
% Needs XeLaTeX:  latexmk -xelatex main.tex
% ============================================================================
\makeatletter
\atldefault{language}{en}
\atldefault{category}{homework}
\newif\ifatlfa
\atlifeq{language}{fa}{\atlfatrue}{\atlfafalse}

% ============================================================================
% atl-packages.tex -- every package the system needs, loaded once, in a safe order.
% Everything here is language-neutral. XePersian (Persian) is loaded afterwards,
% by atl-lang-fa.tex, because it must come after hyperref.
% ============================================================================
\makeatletter
\PassOptionsToPackage{table,dvipsnames,svgnames}{xcolor}
\RequirePackage{etoolbox}
\RequirePackage{fontspec}
\RequirePackage{geometry}
\RequirePackage{setspace}
\RequirePackage{amsmath,amssymb,amsthm,mathtools,bm}
\RequirePackage{siunitx}
\RequirePackage{graphicx}
\RequirePackage{xcolor}
\RequirePackage{array,booktabs,multirow}
\RequirePackage{caption,subcaption}
\RequirePackage{enumitem}
\RequirePackage{needspace}
\RequirePackage{titlesec}
\RequirePackage{fancyhdr}
\RequirePackage{tikz}
\usetikzlibrary{calc,positioning,arrows.meta,shapes.geometric,backgrounds,fit}
\RequirePackage{tcolorbox}
\tcbuselibrary{skins,breakable,listings}
\RequirePackage{listings}
\RequirePackage{hyperref}
\RequirePackage{bookmark}
\RequirePackage{cleveref}
\makeatother

\IfFileExists{preamble/packages.tex}{% بسته‌های اضافه (\usepackage). پیش از لایهٔ زبان بارگذاری می‌شود (XePersian باید آخر بیاید).
% \usepackage{pgfplots}\pgfplotsset{compat=1.18}
% \usepackage{circuitikz}
}{}  % your extra \usepackage lines (before XePersian)
\ifatlfa% ============================================================================
% atl-lang-fa.tex -- Persian (Farsi) language layer, XeLaTeX + XePersian only
%   keys read: fafont = vazirmatn (default) | felmi
%              fascale   (Persian font scale, default depends on the font)
% Fonts are looked up in fonts/ and fall back
% to system fonts, so the document never fails on a missing font.
% ============================================================================
\makeatletter
\RequirePackage{xepersian}

\newcommand{\atl@fa@try}[3]{% regular file, bold file, scale
  % Persian fonts have no italic: map italic to upright so nothing is substituted silently
  \settextfont[Path=fonts/,Scale=#3,BoldFont=#2,ItalicFont=#1,BoldItalicFont=#2]{#1}%
  \setdigitfont[Path=fonts/,Scale=#3,BoldFont=#2,ItalicFont=#1,BoldItalicFont=#2]{#1}%
}
\newcommand{\atl@fa@sys}[1]{% system font family
  \settextfont[Scale=1.0]{#1}\setdigitfont[Scale=1.0]{#1}%
}

\atlifeq{fafont}{felmi}{%
  \atldefault{fascale}{1.1}%
  \IfFileExists{fonts/HM_FElmi.ttf}{\atl@fa@try{HM_FElmi.ttf}{HM_FElmiBd.ttf}{\atlmeta{fascale}}}%
    {\IfFontExistsTF{Amiri}{\atl@fa@sys{Amiri}}{\atl@fa@sys{Noto Naskh Arabic}}}%
  \setlatintextfont[Scale=MatchLowercase]{TeX Gyre Termes}%
}{%
  \atldefault{fascale}{0.95}%
  \IfFileExists{fonts/Vazirmatn-Regular.ttf}{\atl@fa@try{Vazirmatn-Regular.ttf}{Vazirmatn-Bold.ttf}{\atlmeta{fascale}}}%
    {\IfFontExistsTF{Vazirmatn}{\atl@fa@sys{Vazirmatn}}{\atl@fa@sys{Noto Sans Arabic}}}%
  \setlatintextfont[Scale=MatchLowercase]{TeX Gyre Heros}%
}
% Persian sans (headings) is always Vazirmatn when available
\IfFileExists{fonts/Vazirmatn-Regular.ttf}{%
  \setpersiansansfont[Path=fonts/,Scale=0.95,BoldFont=Vazirmatn-Bold.ttf]{Vazirmatn-Regular.ttf}}{}
\IfFontExistsTF{TeX Gyre Heros}{\setsansfont[Scale=MatchLowercase]{TeX Gyre Heros}}{}
\IfFontExistsTF{Latin Modern Mono}{\setmonofont[Scale=MatchLowercase]{Latin Modern Mono}}{}

\AtBeginDocument{%
  \SetSymbolFont{operators}{normal}{OT1}{cmr}{m}{n}%
  \SetSymbolFont{operators}{bold}{OT1}{cmr}{bx}{n}}

% ---- direction hooks (start side of a line = right) -------------------------
\def\atl@start{east}\def\atl@end{west}
\def\atl@L{R}\def\atl@R{L}
% \sffamily is the LATIN sans in XePersian; Persian sans is \persiansffamily
\newcommand{\atl@sf}{\persiansffamily}
% Latin/code runs: left-to-right and no digit conversion (XePersian's inter-character tokens
% would otherwise reorder things like `h:.0e' or `3.14' inside listings)
\newcommand{\atl@ltr}{\setLTR\XeTeXinterchartokenstate=0\resetlatinfont\relax}
\newcommand{\atl@lr}[1]{\lr{#1}}
\newcommand{\atl@numsep}{-}                 % ۲-۱  (Persian academic convention)
\newcommand{\atl@frontnumbering}{harfi}     % آ، ب، پ ...
\newcommand{\atl@appendixnumber}[1]{\harfi{#1}}
\newcommand{\atl@startskip}[1]{\hspace*{#1}}

% ---- labels ----------------------------------------------------------------
\atldefstr{note}{نکته}                 \atldefstr{warning}{هشدار}
\atldefstr{tip}{راهنما}                \atldefstr{result}{نتیجه}
\atldefstr{definition}{تعریف}          \atldefstr{example}{مثال}
\atldefstr{remark}{یادآوری}            \atldefstr{theorem}{قضیه}
\atldefstr{lemma}{لم}                  \atldefstr{proposition}{گزاره}
\atldefstr{corollary}{نتیجه‌ی قضیه}      \atldefstr{proof}{اثبات}
\atldefstr{problem}{مسئله}             \atldefstr{solution}{پاسخ}
\atldefstr{points}{نمره}               \atldefstr{point}{نمره}
\atldefstr{abstract}{چکیده}            \atldefstr{keywords}{واژگان کلیدی}
\atldefstr{contents}{فهرست مطالب}      \atldefstr{listoffigures}{فهرست شکل‌ها}
\atldefstr{listoftables}{فهرست جدول‌ها}
\atldefstr{references}{مراجع}          \atldefstr{appendix}{پیوست}
\atldefstr{chapter}{فصل}               \atldefstr{figure}{شکل}
\atldefstr{table}{جدول}                \atldefstr{page}{صفحه}
\atldefstr{author}{نگارنده}            \atldefstr{authors}{نویسندگان}
\atldefstr{studentid}{شمارهٔ دانشجویی}   \atldefstr{course}{درس}
\atldefstr{coursecode}{کد درس}         \atldefstr{instructor}{استاد}
\atldefstr{supervisor}{استاد راهنما}    \atldefstr{advisor}{استاد مشاور}
\atldefstr{team}{اعضای گروه}           \atldefstr{partners}{همکاران آزمایش}
\atldefstr{experiment}{آزمایش}         \atldefstr{labsection}{گروه آزمایشگاه}
\atldefstr{term}{نیم‌سال}               \atldefstr{date}{تاریخ}
\atldefstr{duedate}{مهلت تحویل}        \atldefstr{submissiondate}{تاریخ تحویل}
\atldefstr{performed}{تاریخ انجام}
\atldefstr{university}{دانشگاه}        \atldefstr{department}{دانشکده}
\atldefstr{type-homework}{تمرین}
\atldefstr{type-lab-report}{گزارش آزمایشگاه}
\atldefstr{type-project-report}{گزارش پروژه}
\atldefstr{type-formal-report}{گزارش}
\atldefstr{type-essay}{مقاله‌ی تحلیلی}
\atldefstr{type-paper}{مقاله}

% cleveref names
\crefname{figure}{شکل}{شکل‌ها}\Crefname{figure}{شکل}{شکل‌ها}
\crefname{table}{جدول}{جدول‌ها}\Crefname{table}{جدول}{جدول‌ها}
\crefname{equation}{رابطه}{رابطه‌ها}\Crefname{equation}{رابطه}{رابطه‌ها}
\crefname{section}{بخش}{بخش‌ها}\Crefname{section}{بخش}{بخش‌ها}
\crefname{chapter}{فصل}{فصل‌ها}\Crefname{chapter}{فصل}{فصل‌ها}
\makeatother
\else% ============================================================================
% atl-lang-en.tex -- English language layer
%   keys read: enfont = termes (default) | pagella | lm
% Provides the same hooks as atl-lang-fa.tex (see the list at the end).
% ============================================================================
\makeatletter
\defaultfontfeatures{Ligatures=TeX}

\newcommand{\atl@setmain}[1]{\IfFontExistsTF{#1}{\setmainfont{#1}}{}}
\atlifeq{enfont}{lm}{}{%
  \atlifeq{enfont}{pagella}{\atl@setmain{TeX Gyre Pagella}}{\atl@setmain{TeX Gyre Termes}}}
\IfFontExistsTF{TeX Gyre Heros}{\setsansfont[Scale=MatchLowercase]{TeX Gyre Heros}}{}
\IfFontExistsTF{Latin Modern Mono}{\setmonofont[Scale=MatchLowercase]{Latin Modern Mono}}{}

% fontspec points the math 'operators' font at the text font, which breaks \mathrm{\Omega}
% and friends (upright Greek); keep upright math symbols in Computer Modern.
\AtBeginDocument{%
  \SetSymbolFont{operators}{normal}{OT1}{cmr}{m}{n}%
  \SetSymbolFont{operators}{bold}{OT1}{cmr}{bx}{n}}

% cleveref names
\crefname{figure}{Figure}{Figures}\Crefname{figure}{Figure}{Figures}
\crefname{table}{Table}{Tables}\Crefname{table}{Table}{Tables}
\crefname{equation}{Equation}{Equations}\Crefname{equation}{Equation}{Equations}
\crefname{section}{Section}{Sections}\Crefname{section}{Section}{Sections}
\crefname{chapter}{Chapter}{Chapters}\Crefname{chapter}{Chapter}{Chapters}

% ---- direction hooks (start side of a line = left) --------------------------
\def\atl@start{west}\def\atl@end{east}      % TikZ / tcolorbox sides
\def\atl@L{L}\def\atl@R{R}                  % fancyhdr: start / end of the line
\newcommand{\atl@sf}{\sffamily}                 % sans family for headings / labels
\newcommand{\atl@ltr}{}                     % force left-to-right (code, bibliography)
\newcommand{\atl@lr}[1]{#1}                 % inline left-to-right run
\newcommand{\atl@numsep}{.}                 % 2.1  (chapter . item)
\newcommand{\atl@frontnumbering}{roman}     % \pagenumbering style for front matter
\newcommand{\atl@appendixnumber}[1]{\Alph{#1}}
\newcommand{\atl@startskip}[1]{\hspace*{#1}}% indent from the start side

% ---- labels ----------------------------------------------------------------
\atldefstr{note}{Note}                 \atldefstr{warning}{Warning}
\atldefstr{tip}{Tip}                   \atldefstr{result}{Result}
\atldefstr{definition}{Definition}     \atldefstr{example}{Example}
\atldefstr{remark}{Remark}             \atldefstr{theorem}{Theorem}
\atldefstr{lemma}{Lemma}               \atldefstr{proposition}{Proposition}
\atldefstr{corollary}{Corollary}       \atldefstr{proof}{Proof}
\atldefstr{problem}{Problem}           \atldefstr{solution}{Solution}
\atldefstr{points}{points}             \atldefstr{point}{point}
\atldefstr{abstract}{Abstract}         \atldefstr{keywords}{Keywords}
\atldefstr{contents}{Contents}         \atldefstr{listoffigures}{List of Figures}
\atldefstr{listoftables}{List of Tables}
\atldefstr{references}{References}     \atldefstr{appendix}{Appendix}
\atldefstr{chapter}{Chapter}           \atldefstr{figure}{Figure}
\atldefstr{table}{Table}               \atldefstr{page}{Page}
\atldefstr{author}{Author}             \atldefstr{authors}{Authors}
\atldefstr{studentid}{Student ID}      \atldefstr{course}{Course}
\atldefstr{coursecode}{Course code}    \atldefstr{instructor}{Instructor}
\atldefstr{supervisor}{Supervisor}     \atldefstr{advisor}{Advisor}
\atldefstr{team}{Team}                 \atldefstr{partners}{Lab partners}
\atldefstr{experiment}{Experiment}     \atldefstr{labsection}{Lab section}
\atldefstr{term}{Term}                 \atldefstr{date}{Date}
\atldefstr{duedate}{Due date}          \atldefstr{submissiondate}{Submitted}
\atldefstr{performed}{Performed on}
\atldefstr{university}{University}     \atldefstr{department}{Department}
\atldefstr{type-homework}{Homework}
\atldefstr{type-lab-report}{Lab Report}
\atldefstr{type-project-report}{Project Report}
\atldefstr{type-formal-report}{Report}
\atldefstr{type-essay}{Essay}
\atldefstr{type-paper}{Paper}
\makeatother
% Hooks every language layer must define:
%   \atl@start \atl@end \atl@L \atl@R \atl@ltr \atl@lr \atl@numsep
%   \atl@sf \atl@frontnumbering \atl@appendixnumber \atl@startskip, and all \atldefstr labels.
\fi
% ============================================================================
% atl-palettes.tex -- colour palettes
%
%   Choose:   \atlset{palette=sunset}   (config.tex)   or   \atlsetpalette{sunset}
%   Define:   \atldefpalette{mine}{primary=003D82, secondary=0066CC, accent=D4AF37,
%                                  light=E6F7FF, tint=EAF2FF}
%   (success, warning, info, result are optional; they fall back to shared defaults.)
%
% Roles used by every other module (never use raw colours in templates):
%   atl-primary     headings, title, table header, links        (dark, readable)
%   atl-secondary   sub-headings, rules                         (mid tone)
%   atl-accent      thin rules, highlights, bars                (decorative)
%   atl-light       box background        atl-tint   page/band tint
%   atl-success atl-warning atl-info atl-result   semantic colours
% Derived: atl-primary-dark, atl-secondary-text (readable on white), atl-accent-text,
%          atl-muted, atl-ink (body text)
% ============================================================================
\makeatletter
\definecolor{atl-default-success}{HTML}{2E7D32}
\definecolor{atl-default-warning}{HTML}{C62828}
\definecolor{atl-default-info}{HTML}{1565C0}
\definecolor{atl-default-result}{HTML}{6A1B9A}

\pgfkeys{/atl/pal/.cd,
  primary/.code={\definecolor{atl-\atl@curpal-primary}{HTML}{#1}},
  secondary/.code={\definecolor{atl-\atl@curpal-secondary}{HTML}{#1}},
  accent/.code={\definecolor{atl-\atl@curpal-accent}{HTML}{#1}},
  light/.code={\definecolor{atl-\atl@curpal-light}{HTML}{#1}},
  tint/.code={\definecolor{atl-\atl@curpal-tint}{HTML}{#1}},
  success/.code={\definecolor{atl-\atl@curpal-success}{HTML}{#1}},
  warning/.code={\definecolor{atl-\atl@curpal-warning}{HTML}{#1}},
  info/.code={\definecolor{atl-\atl@curpal-info}{HTML}{#1}},
  result/.code={\definecolor{atl-\atl@curpal-result}{HTML}{#1}},
}
\newcommand{\atl@palfill}[1]{%
  \ifcsname\string\color@atl-\atl@curpal-#1\endcsname\else
    \colorlet{atl-\atl@curpal-#1}{atl-default-#1}\fi}
\newcommand{\atldefpalette}[2]{%
  \def\atl@curpal{#1}\pgfkeys{/atl/pal/.cd,#2}%
  \forcsvlist{\atl@palfill}{success,warning,info,result}}

% name        primary  secondary accent  light   tint
\atldefpalette{blue}{primary=003D82, secondary=0066CC, accent=D4AF37, light=E6F2FF, tint=EAF2FF}
\atldefpalette{red}{primary=8B0000, secondary=C62828, accent=E0A800, light=FFF0F0, tint=FFE8E8}
\atldefpalette{gray}{primary=2C3E50, secondary=5D6D7E, accent=3498DB, light=ECEFF1, tint=F4F6F7}
\atldefpalette{cyan}{primary=00687A, secondary=00A0B2, accent=FF6F00, light=E1F5FE, tint=E0F2F1}
\atldefpalette{magenta}{primary=7B1FA2, secondary=C2185B, accent=00ACC1, light=F3E5F5, tint=FCE4EC}
\atldefpalette{yellow}{primary=826200, secondary=CD8500, accent=1976D2, light=FFF8E1, tint=FFF3C4}
\atldefpalette{green}{primary=1B5E20, secondary=2E7D32, accent=F9A825, light=E8F5E9, tint=F1F8E9}
\atldefpalette{sunset}{primary=355C7D, secondary=6C5B7B, accent=F67280, light=FAF4F5, tint=FFF1F2}
\atldefpalette{ocean}{primary=2E3A38, secondary=11999E, accent=30E3CA, light=E4F9F5, tint=EEFBF8}
\atldefpalette{berry}{primary=311D3F, secondary=522546, accent=E23E57, light=F8EEF1, tint=FBF3F5}
\atldefpalette{neon}{primary=252A34, secondary=08D9D6, accent=FF2E63, light=EAEAEA, tint=F4F4F4}

\newcommand{\atl@applyrole}[2]{\colorlet{atl-#2}{atl-#1-#2}}
\newcommand{\atlsetpalette}[1]{%
  \ifcsname\string\color@atl-#1-primary\endcsname\else
    \PackageError{atl}{Unknown palette `#1'}{Palettes: blue red gray cyan magenta yellow green sunset ocean berry neon, or define one with \string\atldefpalette.}%
  \fi
  \forcsvlist{\atl@applyrole{#1}}{primary,secondary,accent,light,tint,success,warning,info,result}%
  \colorlet{atl-primary-dark}{atl-primary!70!black}%
  \colorlet{atl-secondary-text}{atl-secondary!78!black}%
  \colorlet{atl-accent-text}{atl-accent!72!black}%
  \colorlet{atl-muted}{black!58}%
  \colorlet{atl-ink}{black!92}%
}
\atldefault{palette}{blue}
\edef\atl@tmp{\atlmeta{palette}}
\expandafter\atlsetpalette\expandafter{\atl@tmp}

\hypersetup{colorlinks=true,linkcolor=atl-primary,urlcolor=atl-secondary-text,citecolor=atl-primary,pdfborder={0 0 0}}
\makeatother

% ============================================================================
% atl-layout.tex -- page geometry, spacing, paragraph style
%   keys: margins   = normal (default) | narrow | formal | wide
%         linespread= 1.15 (default)   (setspace factor; formal reports use 1.5)
%         parstyle  = block (default) | indent
%   "formal" gives the binding side (start of line: left in English, right in
%   Persian) 3 cm and the other side 2.5 cm, as in Iranian thesis rules.
% ============================================================================
\makeatletter
\atldefault{margins}{normal}
\atldefault{linespread}{1.15}
\atldefault{parstyle}{block}

\newcommand{\atl@setmargins}[4]{% top/bottom, start, end, headheight
  \ifdefstring{\atl@start}{west}%
    {\geometry{a4paper,top=#1,bottom=#1,left=#2,right=#3,headheight=#4,headsep=0.8cm,footskip=1.1cm}}%
    {\geometry{a4paper,top=#1,bottom=#1,right=#2,left=#3,headheight=#4,headsep=0.8cm,footskip=1.1cm}}}
\atlifeq{margins}{narrow}{\atl@setmargins{2cm}{2cm}{2cm}{16pt}}{%
\atlifeq{margins}{formal}{\atl@setmargins{2.5cm}{3cm}{2.5cm}{16pt}}{%
\atlifeq{margins}{wide}{\atl@setmargins{3cm}{3.5cm}{3.5cm}{16pt}}{%
\atl@setmargins{2.5cm}{2.5cm}{2.5cm}{16pt}}}}

\edef\atl@tmp{\atlmeta{linespread}}
\expandafter\setstretch\expandafter{\atl@tmp}
\atlifeq{parstyle}{indent}{%
  \setlength{\parindent}{0.71cm}\setlength{\parskip}{0pt}%
}{%
  \setlength{\parindent}{0pt}\setlength{\parskip}{0.6em plus 0.15em minus 0.1em}%
}
\setlength{\emergencystretch}{2em}
\raggedbottom
\setlist{itemsep=0.25em,topsep=0.4em,parsep=0pt}
\makeatother

% ============================================================================
% atl-headings.tex -- sectioning styles, numbering, front/main matter, abstract
%   keys: headfont = sans (default) | serif
%         numbering= flat (default) | section      (article class only; a class with
%                    \chapter always numbers floats/equations per chapter)
%         headsize = normal (default) | compact   (compact: smaller headings, for papers)
%         sectionbreak = off (default) | on       (start every \section on a new page)
% ============================================================================
\makeatletter
\atldefault{headfont}{sans}
\atldefault{numbering}{flat}
\atldefault{headsize}{normal}
\atldefault{sectionbreak}{off}
\atlifeq{headfont}{serif}{\newcommand{\atl@headfont}{\rmfamily}}{\newcommand{\atl@headfont}{\atl@sf}}

% ---- numbering ----------------------------------------------------------------
\@ifundefined{chapter}{%
  \renewcommand{\thesubsection}{\thesection\atl@numsep\arabic{subsection}}%
  \renewcommand{\thesubsubsection}{\thesubsection\atl@numsep\arabic{subsubsection}}%
  \atlifeq{numbering}{section}{%
    \counterwithin{equation}{section}\counterwithin{figure}{section}\counterwithin{table}{section}%
    \renewcommand{\theequation}{\thesection\atl@numsep\arabic{equation}}%
    \renewcommand{\thefigure}{\thesection\atl@numsep\arabic{figure}}%
    \renewcommand{\thetable}{\thesection\atl@numsep\arabic{table}}%
  }{}%
}{%
  \renewcommand{\thesection}{\thechapter\atl@numsep\arabic{section}}%
  \renewcommand{\thesubsection}{\thesection\atl@numsep\arabic{subsection}}%
  \renewcommand{\thesubsubsection}{\thesubsection\atl@numsep\arabic{subsubsection}}%
  \renewcommand{\theequation}{\thechapter\atl@numsep\arabic{equation}}%
  \renewcommand{\thefigure}{\thechapter\atl@numsep\arabic{figure}}%
  \renewcommand{\thetable}{\thechapter\atl@numsep\arabic{table}}%
}

% ---- section styles -----------------------------------------------------------
\titleformat{\section}
  {\atl@headfont\Large\bfseries\color{atl-primary}}
  {\color{atl-accent-text}\thesection}{0.7em}{}
  [\vspace{-0.3ex}{\color{atl-accent}\titlerule[0.9pt]}]
\titleformat{\subsection}
  {\atl@headfont\large\bfseries\color{atl-secondary-text}}
  {\thesubsection}{0.6em}{}
\titleformat{\subsubsection}
  {\atl@headfont\normalsize\bfseries\color{atl-muted}}
  {\thesubsubsection}{0.5em}{}
\titleformat{\paragraph}[runin]
  {\atl@headfont\normalsize\bfseries\color{atl-secondary-text}}
  {}{0pt}{}[.]
\titlespacing*{\section}{0pt}{2.4ex plus 0.6ex minus 0.2ex}{1.2ex plus 0.2ex}
\titlespacing*{\subsection}{0pt}{2ex plus 0.5ex minus 0.2ex}{0.9ex plus 0.2ex}
\titlespacing*{\subsubsection}{0pt}{1.6ex plus 0.4ex}{0.6ex}
\titlespacing*{\paragraph}{0pt}{1.2ex plus 0.3ex}{0.8em}

\atlifeq{headsize}{compact}{%
  \titleformat{\section}
    {\atl@headfont\large\bfseries\color{atl-primary}}
    {\color{atl-accent-text}\thesection}{0.6em}{}
    [\vspace{-0.4ex}{\color{atl-accent}\titlerule[0.7pt]}]
  \titleformat{\subsection}
    {\atl@headfont\normalsize\bfseries\color{atl-secondary-text}}
    {\thesubsection}{0.5em}{}
  \titlespacing*{\section}{0pt}{1.6ex plus 0.4ex}{0.8ex}
  \titlespacing*{\subsection}{0pt}{1.3ex plus 0.3ex}{0.5ex}
}{}
% keep a heading together with the text that follows it
\pretocmd{\section}{\needspace{7\baselineskip}}{}{}
\pretocmd{\subsection}{\needspace{5\baselineskip}}{}{}
\pretocmd{\subsubsection}{\needspace{4\baselineskip}}{}{}
\atlifon{sectionbreak}{\pretocmd{\section}{\clearpage}{}{}}{}

\@ifundefined{chapter}{}{%
  \titleformat{\chapter}[display]
    {\atl@headfont\bfseries\color{atl-primary}}
    {\normalfont\large\color{atl-accent-text}\atlstr{chapter}\ \thechapter}{0.5ex}{\Huge}
    [\vspace{1.2ex}{\color{atl-accent}\titlerule[2.2pt]}]
  \titleformat{name=\chapter,numberless}[display]
    {\atl@headfont\bfseries\color{atl-primary}}
    {}{0pt}{\Huge}
    [\vspace{1.2ex}{\color{atl-accent}\titlerule[2.2pt]}]
  \titlespacing*{\chapter}{0pt}{-4pt}{2.2em}
}

% ---- unnumbered headings that still go to the table of contents ----------------
%   \atlunnumbered{Abstract}   -> chapter* in report class, section* in article
\newcommand{\atlunnumbered}[1]{%
  \@ifundefined{chapter}{%
    \phantomsection\section*{#1}\addcontentsline{toc}{section}{#1}%
  }{%
    \clearpage\phantomsection\chapter*{#1}\addcontentsline{toc}{chapter}{#1}%
    \markboth{#1}{#1}%
  }}

% ---- front / main / back matter (page numbering) ----------------------------------
\newcommand{\atlfrontmatter}{\clearpage\pagenumbering{\atl@frontnumbering}}
\newcommand{\atlmainmatter}{\clearpage\pagenumbering{arabic}}
\newcommand{\atlappendix}{%
  \appendix
  \@ifundefined{chapter}{%
    \renewcommand{\thesection}{\atl@appendixnumber{section}}%
    \renewcommand{\thesubsection}{\thesection\atl@numsep\arabic{subsection}}%
  }{%
    \renewcommand{\thechapter}{\atl@appendixnumber{chapter}}%
    \titleformat{\chapter}[display]
      {\atl@headfont\bfseries\color{atl-primary}}
      {\normalfont\large\color{atl-accent-text}\atlstr{appendix}\ \thechapter}{0.5ex}{\Huge}
      [\vspace{1.2ex}{\color{atl-accent}\titlerule[2.2pt]}]%
  }}

% ---- keywords line (use inside or after an abstract) --------------------------------
\newcommand{\atlkeywords}[1]{\par\smallskip\noindent{\atl@headfont\bfseries\color{atl-primary}\atlstr{keywords}:}\ #1\par}

% ---- abstract -------------------------------------------------------------------
\renewenvironment{abstract}{%
  \begin{tcolorbox}[enhanced,breakable,frame hidden,boxrule=0pt,sharp corners,
      colback=atl-tint,borderline \atl@start={3pt}{0pt}{atl-accent},
      left=10pt,right=10pt,top=6pt,bottom=6pt,before skip=1em,after skip=1.4em]
  {\atl@headfont\bfseries\color{atl-primary}\atlstr{abstract}}\par\smallskip
}{\end{tcolorbox}}
\makeatother

% ============================================================================
% atl-pagestyle.tex -- running header and footer (fancyhdr)
%   keys: headerstart  text at the start side of the header (default: course, then title)
%         headerend    text at the end side (default: current section title)
%         header     = on (default) | off
% In Persian "start" is the right-hand side, so the header mirrors automatically.
% ============================================================================
\makeatletter
\atldefault{header}{on}
\newcommand{\atl@hdrstart}{\atlifset{headerstart}{\atlmeta{headerstart}}{\atlifset{course}{\atlmeta{course}}{\atlmeta{title}}}}
\newcommand{\atl@hdrend}{\atlifset{headerend}{\atlmeta{headerend}}{\nouppercase{\rightmark}}}

\renewcommand{\sectionmark}[1]{\markright{\ifnum\c@secnumdepth>\z@\thesection\ \fi #1}}
\@ifundefined{chapter}{}{\renewcommand{\chaptermark}[1]{\markboth{#1}{}}}

% (an hbox with leaders confuses fancyhdr's height check under XePersian; use \hrule)
\def\atl@headrule{{\color{atl-primary}\hrule\@height\headrulewidth\@width\headwidth\vskip-\headrulewidth}}

\pagestyle{fancy}
\fancyhf{}
\expandafter\fancyhead\expandafter[\atl@L]{\small\atl@headfont\bfseries\color{atl-primary}\atl@hdrstart}
\expandafter\fancyhead\expandafter[\atl@R]{\small\atl@headfont\color{atl-muted}\atl@hdrend}
\fancyfoot[C]{\small\color{atl-muted}\thepage}
\renewcommand{\headrulewidth}{0.6pt}
\renewcommand{\footrulewidth}{0pt}
\let\headrule\atl@headrule
\fancypagestyle{plain}{%
  \fancyhf{}\fancyfoot[C]{\small\color{atl-muted}\thepage}%
  \renewcommand{\headrulewidth}{0pt}\renewcommand{\footrulewidth}{0pt}}
\atlifeq{header}{off}{\pagestyle{plain}}{}
\makeatother

% ============================================================================
% atl-floats.tex -- captions, tables, figures, lists
%   Table captions go above, figure captions below (Iranian and IEEE convention).
%   \atlhead{text}   white bold text for a header cell (use after \rowcolor{atl-primary})
%   \atlrows         zebra striping:  \atlrows before \begin{tabular}
% ============================================================================
\makeatletter
\captionsetup{font=small,labelfont={bf,color=atl-primary},labelsep=colon,
  format=plain,justification=centering,singlelinecheck=true,skip=6pt}
\captionsetup{hypcap=false}
\captionsetup[table]{position=above}
\captionsetup[figure]{position=below}
\renewcommand{\arraystretch}{1.18}
\newcommand{\atlhead}[1]{\textcolor{white}{\bfseries #1}}
\newcommand{\atlrows}{\rowcolors{2}{atl-tint}{white}}

\sisetup{
  per-mode=symbol,
  quantity-product={\,},
  inter-unit-product={\cdot},
  detect-weight=true,detect-family=true,
  mode=match,
  table-align-text-post=false,
}
\graphicspath{{fig/}{figures/}{./}}
\DeclareGraphicsExtensions{.pdf,.png,.jpg,.jpeg}
\makeatother

% ============================================================================
% atl-boxes.tex -- callout boxes (tcolorbox) and homework problem/solution
%
%   \begin{notebox}[Optional title] ... \end{notebox}     (also warningbox, tipbox,
%   resultbox, definitionbox, examplebox)
%   \begin{problem}[Title][points] ... \end{problem}
%   \begin{solution} ... \end{solution}
% The coloured bar sits on the start side of the line (left in English, right in
% Persian). All colours come from the active palette.
% ============================================================================
\makeatletter
\tcbset{
  atl callout/.style={
    enhanced,breakable,frame hidden,boxrule=0pt,sharp corners,
    colback=#1!7!white,
    borderline \atl@start={3.2pt}{0pt}{#1},
    left=10pt,right=10pt,top=5pt,bottom=5pt,
    before skip=0.9em,after skip=0.9em,
    fonttitle=\atl@headfont\bfseries,coltitle=#1!82!black,
    attach title to upper={\par\smallskip},
  },
}
\newcommand{\atl@newcallout}[3]{% env name, colour, string key
  \newtcolorbox{#1}[1][\atlstr{#3}]{atl callout=#2,title={##1}}}
\atl@newcallout{notebox}{atl-info}{note}
\atl@newcallout{warningbox}{atl-warning}{warning}
\atl@newcallout{tipbox}{atl-success}{tip}
\atl@newcallout{resultbox}{atl-result}{result}
\atl@newcallout{definitionbox}{atl-primary}{definition}
\atl@newcallout{examplebox}{atl-secondary-text}{example}

% ---- homework: problem / solution ---------------------------------------------
\newcounter{atlproblem}
\NewDocumentEnvironment{problem}{O{} O{}}{%
  \refstepcounter{atlproblem}%
  \begin{tcolorbox}[enhanced,breakable,sharp corners,boxrule=0.6pt,
    colframe=atl-primary,colback=atl-light,
    left=10pt,right=10pt,top=6pt,bottom=6pt,before skip=1.2em,after skip=0.6em,
    fonttitle=\atl@headfont\bfseries,coltitle=white,colbacktitle=atl-primary,
    toptitle=3pt,bottomtitle=3pt,
    title={\atlstr{problem}~\theatlproblem\ifstrempty{#1}{}{\,\textemdash\,#1}%
      \ifstrempty{#2}{}{\hfill\normalfont\small #2~\atlstr{points}}}]%
}{\end{tcolorbox}}
\NewDocumentEnvironment{solution}{O{}}{%
  \begin{tcolorbox}[enhanced,breakable,sharp corners,boxrule=0pt,frame hidden,
    colback=white,borderline \atl@start={2.4pt}{0pt}{atl-secondary},
    left=12pt,right=8pt,top=4pt,bottom=4pt,before skip=0.2em,after skip=1.2em,
    fonttitle=\atl@headfont\bfseries,coltitle=atl-secondary-text,
    attach title to upper={\par\smallskip},
    title={\atlstr{solution}\ifstrempty{#1}{}{\,\textemdash\,#1}}]%
}{\end{tcolorbox}}
\makeatother

% ============================================================================
% atl-code.tex -- syntax-highlighted code
%   key:   listing = darcula (default) | dark | bright | github
%   switch mid-document:  \atlsetlisting{bright}
%   plain listings:       \begin{lstlisting}[language=Python] ... \end{lstlisting}
%   framed with title:    \begin{codebox}{Python} ... \end{codebox}
%                         \begin{codebox}[title=Fast Fourier transform]{Python}
%   from a file:          \atlinputcode[title=main.c]{C}{../simulation/main.c}
% Languages: Python Matlab C C++ Verilog VHDL Assembler(x86masm) bash ... (any of listings').
% Code is always typeset left-to-right, also inside Persian documents.
% Legacy theme names dark_2026 and bright_2026 are accepted.
% ============================================================================
\makeatletter
% \atldeflisttheme{name}{bg=..,fg=..,kw=..,str=..,cmt=..,num=..,bi=..,deco=..,op=..,br=..}  (HTML hex)
\pgfkeys{/atl/lst/.cd,
  bg/.code={\definecolor{atl-\atl@curlt-bg}{HTML}{#1}},   fg/.code={\definecolor{atl-\atl@curlt-fg}{HTML}{#1}},
  kw/.code={\definecolor{atl-\atl@curlt-kw}{HTML}{#1}},   str/.code={\definecolor{atl-\atl@curlt-str}{HTML}{#1}},
  cmt/.code={\definecolor{atl-\atl@curlt-cmt}{HTML}{#1}}, num/.code={\definecolor{atl-\atl@curlt-num}{HTML}{#1}},
  bi/.code={\definecolor{atl-\atl@curlt-bi}{HTML}{#1}},   deco/.code={\definecolor{atl-\atl@curlt-deco}{HTML}{#1}},
  op/.code={\definecolor{atl-\atl@curlt-op}{HTML}{#1}},   br/.code={\definecolor{atl-\atl@curlt-br}{HTML}{#1}},
}
\newcommand{\atldeflisttheme}[2]{\def\atl@curlt{#1}\pgfkeys{/atl/lst/.cd,#2}}
\atldeflisttheme{darcula}{bg=2B2B2B,fg=A9B7C6,kw=CC7832,str=6A8759,cmt=808080,num=6897BB,bi=8888C6,deco=BBB529,op=A9B7C6,br=A9B7C6}
\atldeflisttheme{dark}{bg=141820,fg=C8D2DC,kw=FF8C32,str=78C878,cmt=646E78,num=64B4FF,bi=B496FF,deco=DCB432,op=C8D2DC,br=FF7878}
\atldeflisttheme{bright}{bg=FAFAFA,fg=282832,kw=B45014,str=148232,cmt=78828C,num=0064C8,bi=783CC8,deco=B47800,op=282832,br=C81E32}
\atldeflisttheme{github}{bg=0D1117,fg=C9D1D9,kw=FF7B72,str=A5D6FF,cmt=8B949E,num=79C0FF,bi=D2A8FF,deco=FFA657,op=C9D1D9,br=C9D1D9}

\newcommand{\atlsetlisting}[1]{%
  \def\atl@lt{#1}%
  \ifstrequal{#1}{dark_2026}{\def\atl@lt{dark}}{}%
  \ifstrequal{#1}{bright_2026}{\def\atl@lt{bright}}{}%
  \ifcsname\string\color@atl-\atl@lt-bg\endcsname\else
    \PackageError{atl}{Unknown listing theme `#1'}{Use darcula, dark, bright or github.}\fi
  \forcsvlist{\atl@lstrole}{bg,fg,kw,str,cmt,num,bi,deco,op,br}}
\newcommand{\atl@lstrole}[1]{\colorlet{atl-lst-#1}{atl-\atl@lt-#1}}

% numbers, operators and brackets get their own colour through literate rules
\lstdefinestyle{atl-lst}{
  basicstyle=\atl@ltr\ttfamily\footnotesize\color{atl-lst-fg},
  identifierstyle=\color{atl-lst-fg},
  keywordstyle=\color{atl-lst-kw},
  keywordstyle=[2]\color{atl-lst-bi},
  keywordstyle=[3]\color{atl-lst-deco},
  stringstyle=\color{atl-lst-str},
  commentstyle=\color{atl-lst-cmt}\itshape,
  numbers=left,numberstyle=\atl@ltr\tiny\color{atl-lst-cmt},numbersep=9pt,
  backgroundcolor=\color{atl-lst-bg},
  frame=single,rulecolor=\color{atl-lst-bg},framerule=0pt,framesep=7pt,
  xleftmargin=\dimexpr2.2em+7pt\relax,framexleftmargin=2.2em,xrightmargin=7pt,framexrightmargin=0pt,
  breaklines=true,breakatwhitespace=false,
  tabsize=4,showspaces=false,showstringspaces=false,upquote=true,
  columns=fullflexible,keepspaces=true,
  literate=
    *{0}{{\textcolor{atl-lst-num}{0}}}{1}
     {1}{{\textcolor{atl-lst-num}{1}}}{1}
     {2}{{\textcolor{atl-lst-num}{2}}}{1}
     {3}{{\textcolor{atl-lst-num}{3}}}{1}
     {4}{{\textcolor{atl-lst-num}{4}}}{1}
     {5}{{\textcolor{atl-lst-num}{5}}}{1}
     {6}{{\textcolor{atl-lst-num}{6}}}{1}
     {7}{{\textcolor{atl-lst-num}{7}}}{1}
     {8}{{\textcolor{atl-lst-num}{8}}}{1}
     {9}{{\textcolor{atl-lst-num}{9}}}{1}
     {=}{{\textcolor{atl-lst-op}{=}}}{1}
     {+}{{\textcolor{atl-lst-op}{+}}}{1}
     {-}{{\textcolor{atl-lst-op}{-}}}{1}
     {*}{{\textcolor{atl-lst-op}{*}}}{1}
     {/}{{\textcolor{atl-lst-op}{/}}}{1}
     {,}{{\textcolor{atl-lst-op}{,}}}{1}
     {:}{{\textcolor{atl-lst-op}{:}}}{1}
     {<}{{\textcolor{atl-lst-op}{<}}}{1}
     {>}{{\textcolor{atl-lst-op}{>}}}{1}
     {!}{{\textcolor{atl-lst-op}{!}}}{1}
     {?}{{\textcolor{atl-lst-op}{?}}}{1}
     {.}{{\textcolor{atl-lst-op}{.}}}{1}
     {;}{{\textcolor{atl-lst-op}{;}}}{1}
     {(}{{\textcolor{atl-lst-br}{(}}}{1}
     {)}{{\textcolor{atl-lst-br}{)}}}{1}
     {[}{{\textcolor{atl-lst-br}{[}}}{1}
     {]}{{\textcolor{atl-lst-br}{]}}}{1}
     {\{}{{\textcolor{atl-lst-br}{\{}}}{1}
     {\}}{{\textcolor{atl-lst-br}{\}}}}{1}
     {==}{{\textcolor{atl-lst-op}{==}}}{2}
     {!=}{{\textcolor{atl-lst-op}{!=}}}{2}
     {<=}{{\textcolor{atl-lst-op}{<=}}}{2}
     {>=}{{\textcolor{atl-lst-op}{>=}}}{2},
}

% a few extra built-ins per language (second keyword class)
\lstdefinelanguage{atlPython}[]{Python}{
  morekeywords=[2]{print,len,range,enumerate,open,zip,map,filter,sorted,sum,min,max,abs,int,float,str,list,dict,set,tuple,isinstance,super,type},
  morekeywords=[3]{self,cls,True,False,None}}
\lstdefinelanguage{atlC}[]{C}{
  morekeywords=[2]{printf,scanf,malloc,calloc,free,memcpy,memset,strlen,uint8_t,uint16_t,uint32_t,int8_t,int16_t,int32_t,size_t,bool},
  morekeywords=[3]{NULL,true,false}}
% friendly names -> our extended languages (plain lstlisting keeps the stock lexers)
\def\atl@l@Python{atlPython}\def\atl@l@python{atlPython}\def\atl@l@C{atlC}
\newcommand{\atl@maplang}[1]{\ifcsname atl@l@#1\endcsname\csname atl@l@#1\endcsname\else #1\fi}
\lstset{style=atl-lst}
\atldefault{listing}{darcula}
\edef\atl@tmp{\atlmeta{listing}}
\expandafter\atlsetlisting\expandafter{\atl@tmp}

% ---- framed code box ----------------------------------------------------------------
\newtcblisting{codebox}[2][]{%
  enhanced,breakable,listing only,
  listing options/.expanded={style=atl-lst,language=\atl@maplang{#2},numbers=none,frame=none,xleftmargin=0pt,framexleftmargin=0pt,backgroundcolor={}},
  colback=atl-lst-bg,colframe=atl-lst-bg!70!atl-lst-fg,boxrule=0.5pt,arc=3pt,
  borderline \atl@start={3pt}{0pt}{atl-accent},
  left=10pt,right=8pt,top=5pt,bottom=5pt,before skip=0.8em,after skip=0.8em,
  fonttitle=\atl@ltr\small\ttfamily,coltitle=atl-lst-fg,colbacktitle=atl-lst-bg!88!atl-lst-fg,
  title={#2},#1}
\newcommand{\atlinputcode}[3][]{%
  \tcbinputlisting{enhanced,breakable,listing only,
    listing file={#3},
    listing options/.expanded={style=atl-lst,language=\atl@maplang{#2},numbers=none,frame=none,xleftmargin=0pt,framexleftmargin=0pt,backgroundcolor={}},
    colback=atl-lst-bg,colframe=atl-lst-bg!70!atl-lst-fg,boxrule=0.5pt,arc=3pt,
    borderline \atl@start={3pt}{0pt}{atl-accent},
    left=10pt,right=8pt,top=5pt,bottom=5pt,before skip=0.8em,after skip=0.8em,
    fonttitle=\atl@ltr\small\ttfamily,coltitle=atl-lst-fg,colbacktitle=atl-lst-bg!88!atl-lst-fg,
    title={#2},#1}}
\makeatother

% ============================================================================
% atl-macros.tex -- text helpers and a small, general engineering/maths notation set
%
% Conventions (IEEE style): variables italic; descriptive subscripts, operators and
% units upright; vectors and matrices bold italic through ONE command each (\vect,
% \mat) so the look can be changed in one place. Units always go through siunitx:
%   \valunit{10}{\kilo\ohm}   \qty{2.2}{\micro\farad}   \num{1.5e-3}
% No LaTeX built-ins are redefined (\H \P \Re \Im \vec \div ...).
% Domain-specific packs (electronics, power, RF, ML ...) belong in preamble/local.tex.
% ============================================================================
\makeatletter
% ---- text helpers ----------------------------------------------------------------
\newcommand{\keyterm}[1]{\textbf{\textcolor{atl-primary}{#1}}}
\newcommand{\alert}[1]{\textbf{\textcolor{atl-warning}{#1}}}
\newcommand{\ok}[1]{\textbf{\textcolor{atl-success}{#1}}}
\newcommand{\highlight}[1]{\colorbox{atl-accent!28}{\strut #1}}
\newcommand{\muted}[1]{\textcolor{atl-muted}{#1}}
\newcommand{\todo}[1]{\textcolor{atl-warning}{\textbf{[TODO: #1]}}}

% ---- notation ---------------------------------------------------------------------
\newcommand{\qtag}[2]{\ensuremath{#1_{\mathrm{#2}}}}              % R_in
\newcommand{\qindex}[2]{\ensuremath{#1_{#2}}}                       % x_n
\newcommand{\qtagindex}[3]{\ensuremath{#1_{\mathrm{#2},#3}}}       % I_ref,k
\newcommand{\vect}[1]{\ensuremath{\bm{#1}}}
\newcommand{\mat}[1]{\ensuremath{\bm{#1}}}
\newcommand{\uvec}[1]{\ensuremath{\hat{\bm{#1}}}}
\newcommand{\phasor}[1]{\ensuremath{\tilde{#1}}}
\newcommand{\conj}[1]{\ensuremath{\overline{#1}}}
\newcommand{\trans}[1]{\ensuremath{{#1}^{\mathsf{T}}}}
\newcommand{\herm}[1]{\ensuremath{{#1}^{\mathsf{H}}}}
\newcommand{\pinv}[1]{\ensuremath{{#1}^{\dagger}}}
\DeclarePairedDelimiter{\abs}{\lvert}{\rvert}
\DeclarePairedDelimiter{\norm}{\lVert}{\rVert}
\DeclarePairedDelimiter{\inner}{\langle}{\rangle}
\newcommand{\ddiff}{\ensuremath{\mathop{}\!\mathrm{d}}}
\newcommand{\deriv}[2]{\ensuremath{\frac{\mathrm{d}#1}{\mathrm{d}#2}}}
\newcommand{\nderiv}[3]{\ensuremath{\frac{\mathrm{d}^{#3}#1}{\mathrm{d}#2^{#3}}}}
\newcommand{\pderiv}[2]{\ensuremath{\frac{\partial#1}{\partial#2}}}
\newcommand{\integral}[4]{\ensuremath{\int_{#1}^{#2}#3\,\mathrm{d}#4}}
\newcommand{\econst}{\ensuremath{\mathrm{e}}}
\newcommand{\junit}{\ensuremath{\mathrm{j}}}
\newcommand{\iunit}{\ensuremath{\mathrm{i}}}
\newcommand{\jw}{\ensuremath{\junit\omega}}
\newcommand{\N}{\ensuremath{\mathbb{N}}}
\newcommand{\Z}{\ensuremath{\mathbb{Z}}}
\newcommand{\Q}{\ensuremath{\mathbb{Q}}}
\newcommand{\R}{\ensuremath{\mathbb{R}}}
\newcommand{\C}{\ensuremath{\mathbb{C}}}
\DeclareMathOperator{\sgn}{sgn}
\DeclareMathOperator{\tr}{tr}
\DeclareMathOperator{\rank}{rank}
\DeclareMathOperator{\diag}{diag}
\DeclareMathOperator{\sinc}{sinc}
\DeclareMathOperator{\rect}{rect}
\DeclareMathOperator{\Expect}{\mathbb{E}}
\DeclareMathOperator{\Var}{Var}
\DeclareMathOperator{\Cov}{Cov}
\DeclareMathOperator{\Prob}{\mathbb{P}}
\DeclareMathOperator*{\argmin}{arg\,min}
\DeclareMathOperator*{\argmax}{arg\,max}

% ---- units (siunitx v3) -------------------------------------------------------------
\newcommand{\valunit}[3][]{\qty[#1]{#2}{#3}}
\DeclareSIUnit{\dBm}{dBm}
\DeclareSIUnit{\var}{var}
\DeclareSIUnit{\rpm}{rpm}
\DeclareSIUnit{\ppm}{ppm}
\DeclareSIUnit{\pu}{p.u.}
\makeatother

% ============================================================================
% atl-theorems.tex -- theorem-like environments (amsthm), coloured by the palette
%   theorem lemma proposition corollary (shared counter)  definition example remark  proof
% Labels come from the language layer; bodies are upright (Persian has no italic).
% ============================================================================
\makeatletter
\newtheoremstyle{atlplain}
  {0.8em}{0.8em}{\normalfont}{0pt}
  {\atl@headfont\bfseries\color{atl-primary}}{.}{0.5em}
  {\thmname{#1}\thmnumber{\,#2}\thmnote{\normalfont\ (#3)}}
\newtheoremstyle{atldef}
  {0.8em}{0.8em}{\normalfont}{0pt}
  {\atl@headfont\bfseries\color{atl-secondary-text}}{.}{0.5em}
  {\thmname{#1}\thmnumber{\,#2}\thmnote{\normalfont\ (#3)}}
\newtheoremstyle{atlremark}
  {0.6em}{0.6em}{\normalfont}{0pt}
  {\atl@headfont\bfseries\color{atl-muted}}{.}{0.5em}
  {\thmname{#1}\thmnumber{\,#2}\thmnote{\normalfont\ (#3)}}

\theoremstyle{atlplain}
\newtheorem{theorem}{\atlstr{theorem}}
\newtheorem{lemma}[theorem]{\atlstr{lemma}}
\newtheorem{proposition}[theorem]{\atlstr{proposition}}
\newtheorem{corollary}[theorem]{\atlstr{corollary}}
\theoremstyle{atldef}
\newtheorem{definition}{\atlstr{definition}}
\newtheorem{example}{\atlstr{example}}
\theoremstyle{atlremark}
\newtheorem{remark}{\atlstr{remark}}
\renewcommand{\proofname}{\atlstr{proof}}
\renewcommand{\qedsymbol}{\textcolor{atl-accent}{$\blacksquare$}}
\crefname{theorem}{\atlstr{theorem}}{\atlstr{theorem}}
\makeatother

% ============================================================================
% atl-titles.tex -- title pages.   Call  \atlmaketitle  right after \begin{document}.
%   key titlepage = block (default)  banner + info grid at the top of page 1
%                   cover             full cover page with palette artwork
%                   classic           centred journal-style title (essay, paper)
%                   none
%   keys used: title subtitle author affiliation studentid course coursecode instructor
%              supervisor advisor team partners experiment labsection term date
%              duedate submissiondate university department logo doctype category
%   infofields = comma list of keys shown in the info grid (labels come from the
%                language layer, so the same list works in English and Persian).
% Empty or unset keys are skipped.
% ============================================================================
\makeatletter
\atldefault{titlepage}{block}
\newcount\atl@cnt
\newlength{\atl@cellw}

\newcommand{\atl@doctype}{\atlifset{doctype}{\atlmeta{doctype}}{\atlstr{type-\atlmeta{category}}}}

\def\atl@amp{&}
\def\atl@endrow{\\[8pt]}
\newcommand{\atl@cellbody}[1]{%
  \global\advance\atl@cnt\@ne
  \parbox[t]{\atl@cellw}{\raggedright
    {\small\atl@headfont\color{atl-muted}\atlstr{#1}}\\[1pt]%
    {\bfseries\atlmeta{#1}}}%
  \ifodd\atl@cnt\expandafter\atl@amp\else\expandafter\atl@endrow\fi}
\newcommand{\atl@cell}[1]{\atlifset{#1}{\atl@cellbody{#1}}{}}
\newcommand{\atlinfogrid}[1]{%
  \begingroup
  \setlength{\atl@cellw}{0.465\linewidth}\global\atl@cnt\z@
  \edef\atl@l{#1}%
  \begin{tabular}{@{}p{\atl@cellw}@{\hspace{0.07\linewidth}}p{\atl@cellw}@{}}
    \expandafter\forcsvlist\expandafter\atl@cell\expandafter{\atl@l}%
    \ifodd\atl@cnt\mbox{}\fi
  \end{tabular}%
  \endgroup}

% ---- block: banner on top of page 1 -------------------------------------------------
\newcommand{\atl@title@block}{%
  \begingroup
  \begin{tcolorbox}[enhanced,sharp corners,frame hidden,colback=atl-primary,
      borderline south={4pt}{0pt}{atl-accent},
      left=15pt,right=15pt,top=13pt,bottom=13pt,before skip=0pt,after skip=0pt]
    {\atl@headfont\small\bfseries\color{white!70!atl-accent}\atl@doctype\par}\vspace{4pt}
    {\atl@headfont\bfseries\fontsize{24}{29}\selectfont\color{white}\atlmeta{title}\par}
    \atlifset{subtitle}{\vspace{4pt}{\atl@headfont\large\color{white!82!atl-primary}\atlmeta{subtitle}\par}}{}
  \end{tcolorbox}
  \vspace{1em}
  \atlinfogrid{\atlmeta{infofields}}\par
  \vspace{0.4em}{\color{atl-accent}\rule{\linewidth}{0.5pt}}\par\vspace{0.6em}
  \endgroup}

% ---- classic: journal-style --------------------------------------------------------
\newcommand{\atl@title@classic}{%
  \begin{center}
    {\atl@headfont\bfseries\LARGE\color{atl-primary}\atlmeta{title}\par}
    \atlifset{subtitle}{\vspace{4pt}{\large\color{atl-muted}\atlmeta{subtitle}\par}}{}
    \vspace{7pt}{\color{atl-accent}\rule{0.22\linewidth}{1.2pt}}\par\vspace{9pt}
    \atlifset{author}{{\large\atlmeta{author}\par}}{}
    \atlifset{affiliation}{\vspace{2pt}{\small\color{atl-muted}\atlmeta{affiliation}\par}}{}
    \atlifset{date}{\vspace{2pt}{\small\atlmeta{date}\par}}{}
  \end{center}\vspace{0.4em}}

% ---- cover: full page ----------------------------------------------------------------
\newcommand{\atl@cover@art}{%
  \begin{tikzpicture}[remember picture,overlay]
    \fill[atl-tint] (current page.south west) rectangle (current page.north east);
    % top ribbons (back to front)
    \path[fill=atl-accent!85!atl-tint] (current page.north west)
      -- ([yshift=-6.6cm]current page.north west)
      to[out=-12,in=188] ([yshift=-2.9cm]current page.north east)
      -- (current page.north east) -- cycle;
    \path[fill=atl-secondary!90!atl-tint] (current page.north west)
      -- ([yshift=-5.7cm]current page.north west)
      to[out=-16,in=190] ([yshift=-2.1cm]current page.north east)
      -- (current page.north east) -- cycle;
    \path[fill=atl-primary] (current page.north west)
      -- ([yshift=-4.7cm]current page.north west)
      to[out=-20,in=192] ([yshift=-1.3cm]current page.north east)
      -- (current page.north east) -- cycle;
    \path[fill=white,fill opacity=0.10] (current page.north west)
      -- ([yshift=-3.1cm]current page.north west)
      to[out=-24,in=195] ([yshift=-0.5cm]current page.north east)
      -- (current page.north east) -- cycle;
    % bottom ribbons
    \path[fill=atl-accent!85!atl-tint] (current page.south east)
      -- ([yshift=2.9cm]current page.south east)
      to[out=172,in=-8] ([yshift=1.2cm]current page.south west)
      -- (current page.south west) -- cycle;
    \path[fill=atl-secondary!90!atl-tint] (current page.south east)
      -- ([yshift=2.2cm]current page.south east)
      to[out=170,in=-10] ([yshift=0.8cm]current page.south west)
      -- (current page.south west) -- cycle;
    \path[fill=atl-primary] (current page.south east)
      -- ([yshift=1.5cm]current page.south east)
      to[out=168,in=-12] ([yshift=0.4cm]current page.south west)
      -- (current page.south west) -- cycle;
  \end{tikzpicture}}

\newcommand{\atl@title@cover}{%
  \begin{titlepage}
    \thispagestyle{empty}\atl@cover@art
    \vspace*{3.4cm}
    \begin{center}
      \atlifset{logo}{\includegraphics[height=2.8cm]{\atlmeta{logo}}\par\vspace{0.7cm}}{}
      \atlifset{university}{{\atl@headfont\Large\bfseries\color{atl-primary}\atlmeta{university}\par}}{}
      \atlifset{department}{\vspace{2pt}{\atl@headfont\large\color{atl-secondary-text}\atlmeta{department}\par}}{}
      \vspace{1.6cm}
      {\atl@headfont\small\bfseries\color{atl-accent-text}\atl@doctype\par}\vspace{8pt}
      {\atl@headfont\bfseries\fontsize{30}{36}\selectfont\color{atl-primary}\atlmeta{title}\par}
      \atlifset{subtitle}{\vspace{8pt}{\atl@headfont\Large\color{atl-muted}\atlmeta{subtitle}\par}}{}
      \vspace{14pt}{\color{atl-accent}\rule{0.28\linewidth}{2pt}}\par
    \end{center}
    \vfill
    \atlinfogrid{\atlmeta{infofields}}
    \vspace{2.6cm}
  \end{titlepage}}

\newcommand{\atlmaketitle}{%
  \atlifeq{titlepage}{cover}{\atl@title@cover}{%
  \atlifeq{titlepage}{classic}{\atl@title@classic}{%
  \atlifeq{titlepage}{none}{}{\atl@title@block}}}}

\AtBeginDocument{%
  \hypersetup{unicode=true,pdftitle={\atlmeta{title}},pdfauthor={\atlmeta{author}}}}
\makeatother

% ============================================================================
% atl-bib.tex -- bibliography (BibTeX + IEEEtran, numeric [1], [2], ...)
%   \atlbibliography{refs}      refs.bib in the project; heading and ToC entry included
%   key bibstyle = IEEEtran (default)
% The list is typeset left-to-right (also in Persian documents) because entries are
% Latin. Cite with \cite{key}. latexmk runs BibTeX automatically.
% ============================================================================
\makeatletter
\atldefault{bibstyle}{IEEEtran}
\renewenvironment{thebibliography}[1]{%
  \atlunnumbered{\atlstr{references}}%
  \begingroup\atl@ltr\small
  \list{\@biblabel{\@arabic\c@enumiv}}{%
    \settowidth\labelwidth{\@biblabel{#1}}%
    \leftmargin\labelwidth \advance\leftmargin\labelsep
    \usecounter{enumiv}\let\p@enumiv\@empty
    \renewcommand\theenumiv{\@arabic\c@enumiv}%
    \itemsep=0.2em}%
  \sloppy\clubpenalty4000\widowpenalty4000\sfcode`\.\@m}
 {\def\@noitemerr{\@latex@warning{Empty `thebibliography' environment}}%
  \endlist\endgroup}
\newcommand{\atlbibliography}[1]{%
  \edef\atl@bs{\atlmeta{bibstyle}}%
  \expandafter\bibliographystyle\expandafter{\atl@bs}%
  \bibliography{#1}}
\makeatother


\makeatletter % (each module ends with \makeatother)
% names LaTeX prints itself (Contents, List of Figures, ...) come from the language layer
\AtBeginDocument{%
  \renewcommand{\contentsname}{\atlstr{contents}}%
  \renewcommand{\listfigurename}{\atlstr{listoffigures}}%
  \renewcommand{\listtablename}{\atlstr{listoftables}}%
  \renewcommand{\figurename}{\atlstr{figure}}%
  \renewcommand{\tablename}{\atlstr{table}}%
  \renewcommand{\appendixname}{\atlstr{appendix}}%
  \@ifundefined{refname}{\newcommand{\refname}{\atlstr{references}}}{\renewcommand{\refname}{\atlstr{references}}}%
  \@ifundefined{bibname}{}{\renewcommand{\bibname}{\atlstr{references}}}%
  \@ifundefined{abstractname}{}{\renewcommand{\abstractname}{\atlstr{abstract}}}}
\hypersetup{linktoc=page}
\makeatother
\makeatother
                         <- this file (also loads preamble/packages.tex)
%     % ماکروها و تنظیمات شخصی شما؛ پس از موتور قالب بارگذاری می‌شود.
% \atlsetpalette{green}                 % تغییر پالت رنگ از همین‌جا
% \atlsetlisting{bright}                % تغییر پوستهٔ کد (darcula, dark, bright, github)
% \newcommand{\A}{\mat{A}}

% ---- ماکروهای این تمرین ----
\definecolor{navy}{HTML}{1F3A68}
\definecolor{accent}{HTML}{B3401E}
\definecolor{lane1}{HTML}{E8F0FE}
\definecolor{lane2}{HTML}{FDF1E3}
\newcommand{\code}[1]{\lr{\texttt{#1}}}                      % نام فراخوانی، ثبات، مسیر
\newcommand{\src}[1]{\par\begingroup\color{gray}\small\noindent\textbf{منبع:}~#1\par\endgroup}
                  your own macros and settings
%     \begin{document} \atlmaketitle ...
% Needs XeLaTeX:  latexmk -xelatex main.tex
% ============================================================================
\makeatletter
\atldefault{language}{en}
\atldefault{category}{homework}
\newif\ifatlfa
\atlifeq{language}{fa}{\atlfatrue}{\atlfafalse}

% ============================================================================
% atl-packages.tex -- every package the system needs, loaded once, in a safe order.
% Everything here is language-neutral. XePersian (Persian) is loaded afterwards,
% by atl-lang-fa.tex, because it must come after hyperref.
% ============================================================================
\makeatletter
\PassOptionsToPackage{table,dvipsnames,svgnames}{xcolor}
\RequirePackage{etoolbox}
\RequirePackage{fontspec}
\RequirePackage{geometry}
\RequirePackage{setspace}
\RequirePackage{amsmath,amssymb,amsthm,mathtools,bm}
\RequirePackage{siunitx}
\RequirePackage{graphicx}
\RequirePackage{xcolor}
\RequirePackage{array,booktabs,multirow}
\RequirePackage{caption,subcaption}
\RequirePackage{enumitem}
\RequirePackage{needspace}
\RequirePackage{titlesec}
\RequirePackage{fancyhdr}
\RequirePackage{tikz}
\usetikzlibrary{calc,positioning,arrows.meta,shapes.geometric,backgrounds,fit}
\RequirePackage{tcolorbox}
\tcbuselibrary{skins,breakable,listings}
\RequirePackage{listings}
\RequirePackage{hyperref}
\RequirePackage{bookmark}
\RequirePackage{cleveref}
\makeatother

\IfFileExists{preamble/packages.tex}{% بسته‌های اضافه (\usepackage). پیش از لایهٔ زبان بارگذاری می‌شود (XePersian باید آخر بیاید).
% \usepackage{pgfplots}\pgfplotsset{compat=1.18}
% \usepackage{circuitikz}
}{}  % your extra \usepackage lines (before XePersian)
\ifatlfa% ============================================================================
% atl-lang-fa.tex -- Persian (Farsi) language layer, XeLaTeX + XePersian only
%   keys read: fafont = vazirmatn (default) | felmi
%              fascale   (Persian font scale, default depends on the font)
% Fonts are looked up in fonts/ and fall back
% to system fonts, so the document never fails on a missing font.
% ============================================================================
\makeatletter
\RequirePackage{xepersian}

\newcommand{\atl@fa@try}[3]{% regular file, bold file, scale
  % Persian fonts have no italic: map italic to upright so nothing is substituted silently
  \settextfont[Path=fonts/,Scale=#3,BoldFont=#2,ItalicFont=#1,BoldItalicFont=#2]{#1}%
  \setdigitfont[Path=fonts/,Scale=#3,BoldFont=#2,ItalicFont=#1,BoldItalicFont=#2]{#1}%
}
\newcommand{\atl@fa@sys}[1]{% system font family
  \settextfont[Scale=1.0]{#1}\setdigitfont[Scale=1.0]{#1}%
}

\atlifeq{fafont}{felmi}{%
  \atldefault{fascale}{1.1}%
  \IfFileExists{fonts/HM_FElmi.ttf}{\atl@fa@try{HM_FElmi.ttf}{HM_FElmiBd.ttf}{\atlmeta{fascale}}}%
    {\IfFontExistsTF{Amiri}{\atl@fa@sys{Amiri}}{\atl@fa@sys{Noto Naskh Arabic}}}%
  \setlatintextfont[Scale=MatchLowercase]{TeX Gyre Termes}%
}{%
  \atldefault{fascale}{0.95}%
  \IfFileExists{fonts/Vazirmatn-Regular.ttf}{\atl@fa@try{Vazirmatn-Regular.ttf}{Vazirmatn-Bold.ttf}{\atlmeta{fascale}}}%
    {\IfFontExistsTF{Vazirmatn}{\atl@fa@sys{Vazirmatn}}{\atl@fa@sys{Noto Sans Arabic}}}%
  \setlatintextfont[Scale=MatchLowercase]{TeX Gyre Heros}%
}
% Persian sans (headings) is always Vazirmatn when available
\IfFileExists{fonts/Vazirmatn-Regular.ttf}{%
  \setpersiansansfont[Path=fonts/,Scale=0.95,BoldFont=Vazirmatn-Bold.ttf]{Vazirmatn-Regular.ttf}}{}
\IfFontExistsTF{TeX Gyre Heros}{\setsansfont[Scale=MatchLowercase]{TeX Gyre Heros}}{}
\IfFontExistsTF{Latin Modern Mono}{\setmonofont[Scale=MatchLowercase]{Latin Modern Mono}}{}

\AtBeginDocument{%
  \SetSymbolFont{operators}{normal}{OT1}{cmr}{m}{n}%
  \SetSymbolFont{operators}{bold}{OT1}{cmr}{bx}{n}}

% ---- direction hooks (start side of a line = right) -------------------------
\def\atl@start{east}\def\atl@end{west}
\def\atl@L{R}\def\atl@R{L}
% \sffamily is the LATIN sans in XePersian; Persian sans is \persiansffamily
\newcommand{\atl@sf}{\persiansffamily}
% Latin/code runs: left-to-right and no digit conversion (XePersian's inter-character tokens
% would otherwise reorder things like `h:.0e' or `3.14' inside listings)
\newcommand{\atl@ltr}{\setLTR\XeTeXinterchartokenstate=0\resetlatinfont\relax}
\newcommand{\atl@lr}[1]{\lr{#1}}
\newcommand{\atl@numsep}{-}                 % ۲-۱  (Persian academic convention)
\newcommand{\atl@frontnumbering}{harfi}     % آ، ب، پ ...
\newcommand{\atl@appendixnumber}[1]{\harfi{#1}}
\newcommand{\atl@startskip}[1]{\hspace*{#1}}

% ---- labels ----------------------------------------------------------------
\atldefstr{note}{نکته}                 \atldefstr{warning}{هشدار}
\atldefstr{tip}{راهنما}                \atldefstr{result}{نتیجه}
\atldefstr{definition}{تعریف}          \atldefstr{example}{مثال}
\atldefstr{remark}{یادآوری}            \atldefstr{theorem}{قضیه}
\atldefstr{lemma}{لم}                  \atldefstr{proposition}{گزاره}
\atldefstr{corollary}{نتیجه‌ی قضیه}      \atldefstr{proof}{اثبات}
\atldefstr{problem}{مسئله}             \atldefstr{solution}{پاسخ}
\atldefstr{points}{نمره}               \atldefstr{point}{نمره}
\atldefstr{abstract}{چکیده}            \atldefstr{keywords}{واژگان کلیدی}
\atldefstr{contents}{فهرست مطالب}      \atldefstr{listoffigures}{فهرست شکل‌ها}
\atldefstr{listoftables}{فهرست جدول‌ها}
\atldefstr{references}{مراجع}          \atldefstr{appendix}{پیوست}
\atldefstr{chapter}{فصل}               \atldefstr{figure}{شکل}
\atldefstr{table}{جدول}                \atldefstr{page}{صفحه}
\atldefstr{author}{نگارنده}            \atldefstr{authors}{نویسندگان}
\atldefstr{studentid}{شمارهٔ دانشجویی}   \atldefstr{course}{درس}
\atldefstr{coursecode}{کد درس}         \atldefstr{instructor}{استاد}
\atldefstr{supervisor}{استاد راهنما}    \atldefstr{advisor}{استاد مشاور}
\atldefstr{team}{اعضای گروه}           \atldefstr{partners}{همکاران آزمایش}
\atldefstr{experiment}{آزمایش}         \atldefstr{labsection}{گروه آزمایشگاه}
\atldefstr{term}{نیم‌سال}               \atldefstr{date}{تاریخ}
\atldefstr{duedate}{مهلت تحویل}        \atldefstr{submissiondate}{تاریخ تحویل}
\atldefstr{performed}{تاریخ انجام}
\atldefstr{university}{دانشگاه}        \atldefstr{department}{دانشکده}
\atldefstr{type-homework}{تمرین}
\atldefstr{type-lab-report}{گزارش آزمایشگاه}
\atldefstr{type-project-report}{گزارش پروژه}
\atldefstr{type-formal-report}{گزارش}
\atldefstr{type-essay}{مقاله‌ی تحلیلی}
\atldefstr{type-paper}{مقاله}

% cleveref names
\crefname{figure}{شکل}{شکل‌ها}\Crefname{figure}{شکل}{شکل‌ها}
\crefname{table}{جدول}{جدول‌ها}\Crefname{table}{جدول}{جدول‌ها}
\crefname{equation}{رابطه}{رابطه‌ها}\Crefname{equation}{رابطه}{رابطه‌ها}
\crefname{section}{بخش}{بخش‌ها}\Crefname{section}{بخش}{بخش‌ها}
\crefname{chapter}{فصل}{فصل‌ها}\Crefname{chapter}{فصل}{فصل‌ها}
\makeatother
\else% ============================================================================
% atl-lang-en.tex -- English language layer
%   keys read: enfont = termes (default) | pagella | lm
% Provides the same hooks as atl-lang-fa.tex (see the list at the end).
% ============================================================================
\makeatletter
\defaultfontfeatures{Ligatures=TeX}

\newcommand{\atl@setmain}[1]{\IfFontExistsTF{#1}{\setmainfont{#1}}{}}
\atlifeq{enfont}{lm}{}{%
  \atlifeq{enfont}{pagella}{\atl@setmain{TeX Gyre Pagella}}{\atl@setmain{TeX Gyre Termes}}}
\IfFontExistsTF{TeX Gyre Heros}{\setsansfont[Scale=MatchLowercase]{TeX Gyre Heros}}{}
\IfFontExistsTF{Latin Modern Mono}{\setmonofont[Scale=MatchLowercase]{Latin Modern Mono}}{}

% fontspec points the math 'operators' font at the text font, which breaks \mathrm{\Omega}
% and friends (upright Greek); keep upright math symbols in Computer Modern.
\AtBeginDocument{%
  \SetSymbolFont{operators}{normal}{OT1}{cmr}{m}{n}%
  \SetSymbolFont{operators}{bold}{OT1}{cmr}{bx}{n}}

% cleveref names
\crefname{figure}{Figure}{Figures}\Crefname{figure}{Figure}{Figures}
\crefname{table}{Table}{Tables}\Crefname{table}{Table}{Tables}
\crefname{equation}{Equation}{Equations}\Crefname{equation}{Equation}{Equations}
\crefname{section}{Section}{Sections}\Crefname{section}{Section}{Sections}
\crefname{chapter}{Chapter}{Chapters}\Crefname{chapter}{Chapter}{Chapters}

% ---- direction hooks (start side of a line = left) --------------------------
\def\atl@start{west}\def\atl@end{east}      % TikZ / tcolorbox sides
\def\atl@L{L}\def\atl@R{R}                  % fancyhdr: start / end of the line
\newcommand{\atl@sf}{\sffamily}                 % sans family for headings / labels
\newcommand{\atl@ltr}{}                     % force left-to-right (code, bibliography)
\newcommand{\atl@lr}[1]{#1}                 % inline left-to-right run
\newcommand{\atl@numsep}{.}                 % 2.1  (chapter . item)
\newcommand{\atl@frontnumbering}{roman}     % \pagenumbering style for front matter
\newcommand{\atl@appendixnumber}[1]{\Alph{#1}}
\newcommand{\atl@startskip}[1]{\hspace*{#1}}% indent from the start side

% ---- labels ----------------------------------------------------------------
\atldefstr{note}{Note}                 \atldefstr{warning}{Warning}
\atldefstr{tip}{Tip}                   \atldefstr{result}{Result}
\atldefstr{definition}{Definition}     \atldefstr{example}{Example}
\atldefstr{remark}{Remark}             \atldefstr{theorem}{Theorem}
\atldefstr{lemma}{Lemma}               \atldefstr{proposition}{Proposition}
\atldefstr{corollary}{Corollary}       \atldefstr{proof}{Proof}
\atldefstr{problem}{Problem}           \atldefstr{solution}{Solution}
\atldefstr{points}{points}             \atldefstr{point}{point}
\atldefstr{abstract}{Abstract}         \atldefstr{keywords}{Keywords}
\atldefstr{contents}{Contents}         \atldefstr{listoffigures}{List of Figures}
\atldefstr{listoftables}{List of Tables}
\atldefstr{references}{References}     \atldefstr{appendix}{Appendix}
\atldefstr{chapter}{Chapter}           \atldefstr{figure}{Figure}
\atldefstr{table}{Table}               \atldefstr{page}{Page}
\atldefstr{author}{Author}             \atldefstr{authors}{Authors}
\atldefstr{studentid}{Student ID}      \atldefstr{course}{Course}
\atldefstr{coursecode}{Course code}    \atldefstr{instructor}{Instructor}
\atldefstr{supervisor}{Supervisor}     \atldefstr{advisor}{Advisor}
\atldefstr{team}{Team}                 \atldefstr{partners}{Lab partners}
\atldefstr{experiment}{Experiment}     \atldefstr{labsection}{Lab section}
\atldefstr{term}{Term}                 \atldefstr{date}{Date}
\atldefstr{duedate}{Due date}          \atldefstr{submissiondate}{Submitted}
\atldefstr{performed}{Performed on}
\atldefstr{university}{University}     \atldefstr{department}{Department}
\atldefstr{type-homework}{Homework}
\atldefstr{type-lab-report}{Lab Report}
\atldefstr{type-project-report}{Project Report}
\atldefstr{type-formal-report}{Report}
\atldefstr{type-essay}{Essay}
\atldefstr{type-paper}{Paper}
\makeatother
% Hooks every language layer must define:
%   \atl@start \atl@end \atl@L \atl@R \atl@ltr \atl@lr \atl@numsep
%   \atl@sf \atl@frontnumbering \atl@appendixnumber \atl@startskip, and all \atldefstr labels.
\fi
% ============================================================================
% atl-palettes.tex -- colour palettes
%
%   Choose:   \atlset{palette=sunset}   (config.tex)   or   \atlsetpalette{sunset}
%   Define:   \atldefpalette{mine}{primary=003D82, secondary=0066CC, accent=D4AF37,
%                                  light=E6F7FF, tint=EAF2FF}
%   (success, warning, info, result are optional; they fall back to shared defaults.)
%
% Roles used by every other module (never use raw colours in templates):
%   atl-primary     headings, title, table header, links        (dark, readable)
%   atl-secondary   sub-headings, rules                         (mid tone)
%   atl-accent      thin rules, highlights, bars                (decorative)
%   atl-light       box background        atl-tint   page/band tint
%   atl-success atl-warning atl-info atl-result   semantic colours
% Derived: atl-primary-dark, atl-secondary-text (readable on white), atl-accent-text,
%          atl-muted, atl-ink (body text)
% ============================================================================
\makeatletter
\definecolor{atl-default-success}{HTML}{2E7D32}
\definecolor{atl-default-warning}{HTML}{C62828}
\definecolor{atl-default-info}{HTML}{1565C0}
\definecolor{atl-default-result}{HTML}{6A1B9A}

\pgfkeys{/atl/pal/.cd,
  primary/.code={\definecolor{atl-\atl@curpal-primary}{HTML}{#1}},
  secondary/.code={\definecolor{atl-\atl@curpal-secondary}{HTML}{#1}},
  accent/.code={\definecolor{atl-\atl@curpal-accent}{HTML}{#1}},
  light/.code={\definecolor{atl-\atl@curpal-light}{HTML}{#1}},
  tint/.code={\definecolor{atl-\atl@curpal-tint}{HTML}{#1}},
  success/.code={\definecolor{atl-\atl@curpal-success}{HTML}{#1}},
  warning/.code={\definecolor{atl-\atl@curpal-warning}{HTML}{#1}},
  info/.code={\definecolor{atl-\atl@curpal-info}{HTML}{#1}},
  result/.code={\definecolor{atl-\atl@curpal-result}{HTML}{#1}},
}
\newcommand{\atl@palfill}[1]{%
  \ifcsname\string\color@atl-\atl@curpal-#1\endcsname\else
    \colorlet{atl-\atl@curpal-#1}{atl-default-#1}\fi}
\newcommand{\atldefpalette}[2]{%
  \def\atl@curpal{#1}\pgfkeys{/atl/pal/.cd,#2}%
  \forcsvlist{\atl@palfill}{success,warning,info,result}}

% name        primary  secondary accent  light   tint
\atldefpalette{blue}{primary=003D82, secondary=0066CC, accent=D4AF37, light=E6F2FF, tint=EAF2FF}
\atldefpalette{red}{primary=8B0000, secondary=C62828, accent=E0A800, light=FFF0F0, tint=FFE8E8}
\atldefpalette{gray}{primary=2C3E50, secondary=5D6D7E, accent=3498DB, light=ECEFF1, tint=F4F6F7}
\atldefpalette{cyan}{primary=00687A, secondary=00A0B2, accent=FF6F00, light=E1F5FE, tint=E0F2F1}
\atldefpalette{magenta}{primary=7B1FA2, secondary=C2185B, accent=00ACC1, light=F3E5F5, tint=FCE4EC}
\atldefpalette{yellow}{primary=826200, secondary=CD8500, accent=1976D2, light=FFF8E1, tint=FFF3C4}
\atldefpalette{green}{primary=1B5E20, secondary=2E7D32, accent=F9A825, light=E8F5E9, tint=F1F8E9}
\atldefpalette{sunset}{primary=355C7D, secondary=6C5B7B, accent=F67280, light=FAF4F5, tint=FFF1F2}
\atldefpalette{ocean}{primary=2E3A38, secondary=11999E, accent=30E3CA, light=E4F9F5, tint=EEFBF8}
\atldefpalette{berry}{primary=311D3F, secondary=522546, accent=E23E57, light=F8EEF1, tint=FBF3F5}
\atldefpalette{neon}{primary=252A34, secondary=08D9D6, accent=FF2E63, light=EAEAEA, tint=F4F4F4}

\newcommand{\atl@applyrole}[2]{\colorlet{atl-#2}{atl-#1-#2}}
\newcommand{\atlsetpalette}[1]{%
  \ifcsname\string\color@atl-#1-primary\endcsname\else
    \PackageError{atl}{Unknown palette `#1'}{Palettes: blue red gray cyan magenta yellow green sunset ocean berry neon, or define one with \string\atldefpalette.}%
  \fi
  \forcsvlist{\atl@applyrole{#1}}{primary,secondary,accent,light,tint,success,warning,info,result}%
  \colorlet{atl-primary-dark}{atl-primary!70!black}%
  \colorlet{atl-secondary-text}{atl-secondary!78!black}%
  \colorlet{atl-accent-text}{atl-accent!72!black}%
  \colorlet{atl-muted}{black!58}%
  \colorlet{atl-ink}{black!92}%
}
\atldefault{palette}{blue}
\edef\atl@tmp{\atlmeta{palette}}
\expandafter\atlsetpalette\expandafter{\atl@tmp}

\hypersetup{colorlinks=true,linkcolor=atl-primary,urlcolor=atl-secondary-text,citecolor=atl-primary,pdfborder={0 0 0}}
\makeatother

% ============================================================================
% atl-layout.tex -- page geometry, spacing, paragraph style
%   keys: margins   = normal (default) | narrow | formal | wide
%         linespread= 1.15 (default)   (setspace factor; formal reports use 1.5)
%         parstyle  = block (default) | indent
%   "formal" gives the binding side (start of line: left in English, right in
%   Persian) 3 cm and the other side 2.5 cm, as in Iranian thesis rules.
% ============================================================================
\makeatletter
\atldefault{margins}{normal}
\atldefault{linespread}{1.15}
\atldefault{parstyle}{block}

\newcommand{\atl@setmargins}[4]{% top/bottom, start, end, headheight
  \ifdefstring{\atl@start}{west}%
    {\geometry{a4paper,top=#1,bottom=#1,left=#2,right=#3,headheight=#4,headsep=0.8cm,footskip=1.1cm}}%
    {\geometry{a4paper,top=#1,bottom=#1,right=#2,left=#3,headheight=#4,headsep=0.8cm,footskip=1.1cm}}}
\atlifeq{margins}{narrow}{\atl@setmargins{2cm}{2cm}{2cm}{16pt}}{%
\atlifeq{margins}{formal}{\atl@setmargins{2.5cm}{3cm}{2.5cm}{16pt}}{%
\atlifeq{margins}{wide}{\atl@setmargins{3cm}{3.5cm}{3.5cm}{16pt}}{%
\atl@setmargins{2.5cm}{2.5cm}{2.5cm}{16pt}}}}

\edef\atl@tmp{\atlmeta{linespread}}
\expandafter\setstretch\expandafter{\atl@tmp}
\atlifeq{parstyle}{indent}{%
  \setlength{\parindent}{0.71cm}\setlength{\parskip}{0pt}%
}{%
  \setlength{\parindent}{0pt}\setlength{\parskip}{0.6em plus 0.15em minus 0.1em}%
}
\setlength{\emergencystretch}{2em}
\raggedbottom
\setlist{itemsep=0.25em,topsep=0.4em,parsep=0pt}
\makeatother

% ============================================================================
% atl-headings.tex -- sectioning styles, numbering, front/main matter, abstract
%   keys: headfont = sans (default) | serif
%         numbering= flat (default) | section      (article class only; a class with
%                    \chapter always numbers floats/equations per chapter)
%         headsize = normal (default) | compact   (compact: smaller headings, for papers)
%         sectionbreak = off (default) | on       (start every \section on a new page)
% ============================================================================
\makeatletter
\atldefault{headfont}{sans}
\atldefault{numbering}{flat}
\atldefault{headsize}{normal}
\atldefault{sectionbreak}{off}
\atlifeq{headfont}{serif}{\newcommand{\atl@headfont}{\rmfamily}}{\newcommand{\atl@headfont}{\atl@sf}}

% ---- numbering ----------------------------------------------------------------
\@ifundefined{chapter}{%
  \renewcommand{\thesubsection}{\thesection\atl@numsep\arabic{subsection}}%
  \renewcommand{\thesubsubsection}{\thesubsection\atl@numsep\arabic{subsubsection}}%
  \atlifeq{numbering}{section}{%
    \counterwithin{equation}{section}\counterwithin{figure}{section}\counterwithin{table}{section}%
    \renewcommand{\theequation}{\thesection\atl@numsep\arabic{equation}}%
    \renewcommand{\thefigure}{\thesection\atl@numsep\arabic{figure}}%
    \renewcommand{\thetable}{\thesection\atl@numsep\arabic{table}}%
  }{}%
}{%
  \renewcommand{\thesection}{\thechapter\atl@numsep\arabic{section}}%
  \renewcommand{\thesubsection}{\thesection\atl@numsep\arabic{subsection}}%
  \renewcommand{\thesubsubsection}{\thesubsection\atl@numsep\arabic{subsubsection}}%
  \renewcommand{\theequation}{\thechapter\atl@numsep\arabic{equation}}%
  \renewcommand{\thefigure}{\thechapter\atl@numsep\arabic{figure}}%
  \renewcommand{\thetable}{\thechapter\atl@numsep\arabic{table}}%
}

% ---- section styles -----------------------------------------------------------
\titleformat{\section}
  {\atl@headfont\Large\bfseries\color{atl-primary}}
  {\color{atl-accent-text}\thesection}{0.7em}{}
  [\vspace{-0.3ex}{\color{atl-accent}\titlerule[0.9pt]}]
\titleformat{\subsection}
  {\atl@headfont\large\bfseries\color{atl-secondary-text}}
  {\thesubsection}{0.6em}{}
\titleformat{\subsubsection}
  {\atl@headfont\normalsize\bfseries\color{atl-muted}}
  {\thesubsubsection}{0.5em}{}
\titleformat{\paragraph}[runin]
  {\atl@headfont\normalsize\bfseries\color{atl-secondary-text}}
  {}{0pt}{}[.]
\titlespacing*{\section}{0pt}{2.4ex plus 0.6ex minus 0.2ex}{1.2ex plus 0.2ex}
\titlespacing*{\subsection}{0pt}{2ex plus 0.5ex minus 0.2ex}{0.9ex plus 0.2ex}
\titlespacing*{\subsubsection}{0pt}{1.6ex plus 0.4ex}{0.6ex}
\titlespacing*{\paragraph}{0pt}{1.2ex plus 0.3ex}{0.8em}

\atlifeq{headsize}{compact}{%
  \titleformat{\section}
    {\atl@headfont\large\bfseries\color{atl-primary}}
    {\color{atl-accent-text}\thesection}{0.6em}{}
    [\vspace{-0.4ex}{\color{atl-accent}\titlerule[0.7pt]}]
  \titleformat{\subsection}
    {\atl@headfont\normalsize\bfseries\color{atl-secondary-text}}
    {\thesubsection}{0.5em}{}
  \titlespacing*{\section}{0pt}{1.6ex plus 0.4ex}{0.8ex}
  \titlespacing*{\subsection}{0pt}{1.3ex plus 0.3ex}{0.5ex}
}{}
% keep a heading together with the text that follows it
\pretocmd{\section}{\needspace{7\baselineskip}}{}{}
\pretocmd{\subsection}{\needspace{5\baselineskip}}{}{}
\pretocmd{\subsubsection}{\needspace{4\baselineskip}}{}{}
\atlifon{sectionbreak}{\pretocmd{\section}{\clearpage}{}{}}{}

\@ifundefined{chapter}{}{%
  \titleformat{\chapter}[display]
    {\atl@headfont\bfseries\color{atl-primary}}
    {\normalfont\large\color{atl-accent-text}\atlstr{chapter}\ \thechapter}{0.5ex}{\Huge}
    [\vspace{1.2ex}{\color{atl-accent}\titlerule[2.2pt]}]
  \titleformat{name=\chapter,numberless}[display]
    {\atl@headfont\bfseries\color{atl-primary}}
    {}{0pt}{\Huge}
    [\vspace{1.2ex}{\color{atl-accent}\titlerule[2.2pt]}]
  \titlespacing*{\chapter}{0pt}{-4pt}{2.2em}
}

% ---- unnumbered headings that still go to the table of contents ----------------
%   \atlunnumbered{Abstract}   -> chapter* in report class, section* in article
\newcommand{\atlunnumbered}[1]{%
  \@ifundefined{chapter}{%
    \phantomsection\section*{#1}\addcontentsline{toc}{section}{#1}%
  }{%
    \clearpage\phantomsection\chapter*{#1}\addcontentsline{toc}{chapter}{#1}%
    \markboth{#1}{#1}%
  }}

% ---- front / main / back matter (page numbering) ----------------------------------
\newcommand{\atlfrontmatter}{\clearpage\pagenumbering{\atl@frontnumbering}}
\newcommand{\atlmainmatter}{\clearpage\pagenumbering{arabic}}
\newcommand{\atlappendix}{%
  \appendix
  \@ifundefined{chapter}{%
    \renewcommand{\thesection}{\atl@appendixnumber{section}}%
    \renewcommand{\thesubsection}{\thesection\atl@numsep\arabic{subsection}}%
  }{%
    \renewcommand{\thechapter}{\atl@appendixnumber{chapter}}%
    \titleformat{\chapter}[display]
      {\atl@headfont\bfseries\color{atl-primary}}
      {\normalfont\large\color{atl-accent-text}\atlstr{appendix}\ \thechapter}{0.5ex}{\Huge}
      [\vspace{1.2ex}{\color{atl-accent}\titlerule[2.2pt]}]%
  }}

% ---- keywords line (use inside or after an abstract) --------------------------------
\newcommand{\atlkeywords}[1]{\par\smallskip\noindent{\atl@headfont\bfseries\color{atl-primary}\atlstr{keywords}:}\ #1\par}

% ---- abstract -------------------------------------------------------------------
\renewenvironment{abstract}{%
  \begin{tcolorbox}[enhanced,breakable,frame hidden,boxrule=0pt,sharp corners,
      colback=atl-tint,borderline \atl@start={3pt}{0pt}{atl-accent},
      left=10pt,right=10pt,top=6pt,bottom=6pt,before skip=1em,after skip=1.4em]
  {\atl@headfont\bfseries\color{atl-primary}\atlstr{abstract}}\par\smallskip
}{\end{tcolorbox}}
\makeatother

% ============================================================================
% atl-pagestyle.tex -- running header and footer (fancyhdr)
%   keys: headerstart  text at the start side of the header (default: course, then title)
%         headerend    text at the end side (default: current section title)
%         header     = on (default) | off
% In Persian "start" is the right-hand side, so the header mirrors automatically.
% ============================================================================
\makeatletter
\atldefault{header}{on}
\newcommand{\atl@hdrstart}{\atlifset{headerstart}{\atlmeta{headerstart}}{\atlifset{course}{\atlmeta{course}}{\atlmeta{title}}}}
\newcommand{\atl@hdrend}{\atlifset{headerend}{\atlmeta{headerend}}{\nouppercase{\rightmark}}}

\renewcommand{\sectionmark}[1]{\markright{\ifnum\c@secnumdepth>\z@\thesection\ \fi #1}}
\@ifundefined{chapter}{}{\renewcommand{\chaptermark}[1]{\markboth{#1}{}}}

% (an hbox with leaders confuses fancyhdr's height check under XePersian; use \hrule)
\def\atl@headrule{{\color{atl-primary}\hrule\@height\headrulewidth\@width\headwidth\vskip-\headrulewidth}}

\pagestyle{fancy}
\fancyhf{}
\expandafter\fancyhead\expandafter[\atl@L]{\small\atl@headfont\bfseries\color{atl-primary}\atl@hdrstart}
\expandafter\fancyhead\expandafter[\atl@R]{\small\atl@headfont\color{atl-muted}\atl@hdrend}
\fancyfoot[C]{\small\color{atl-muted}\thepage}
\renewcommand{\headrulewidth}{0.6pt}
\renewcommand{\footrulewidth}{0pt}
\let\headrule\atl@headrule
\fancypagestyle{plain}{%
  \fancyhf{}\fancyfoot[C]{\small\color{atl-muted}\thepage}%
  \renewcommand{\headrulewidth}{0pt}\renewcommand{\footrulewidth}{0pt}}
\atlifeq{header}{off}{\pagestyle{plain}}{}
\makeatother

% ============================================================================
% atl-floats.tex -- captions, tables, figures, lists
%   Table captions go above, figure captions below (Iranian and IEEE convention).
%   \atlhead{text}   white bold text for a header cell (use after \rowcolor{atl-primary})
%   \atlrows         zebra striping:  \atlrows before \begin{tabular}
% ============================================================================
\makeatletter
\captionsetup{font=small,labelfont={bf,color=atl-primary},labelsep=colon,
  format=plain,justification=centering,singlelinecheck=true,skip=6pt}
\captionsetup{hypcap=false}
\captionsetup[table]{position=above}
\captionsetup[figure]{position=below}
\renewcommand{\arraystretch}{1.18}
\newcommand{\atlhead}[1]{\textcolor{white}{\bfseries #1}}
\newcommand{\atlrows}{\rowcolors{2}{atl-tint}{white}}

\sisetup{
  per-mode=symbol,
  quantity-product={\,},
  inter-unit-product={\cdot},
  detect-weight=true,detect-family=true,
  mode=match,
  table-align-text-post=false,
}
\graphicspath{{fig/}{figures/}{./}}
\DeclareGraphicsExtensions{.pdf,.png,.jpg,.jpeg}
\makeatother

% ============================================================================
% atl-boxes.tex -- callout boxes (tcolorbox) and homework problem/solution
%
%   \begin{notebox}[Optional title] ... \end{notebox}     (also warningbox, tipbox,
%   resultbox, definitionbox, examplebox)
%   \begin{problem}[Title][points] ... \end{problem}
%   \begin{solution} ... \end{solution}
% The coloured bar sits on the start side of the line (left in English, right in
% Persian). All colours come from the active palette.
% ============================================================================
\makeatletter
\tcbset{
  atl callout/.style={
    enhanced,breakable,frame hidden,boxrule=0pt,sharp corners,
    colback=#1!7!white,
    borderline \atl@start={3.2pt}{0pt}{#1},
    left=10pt,right=10pt,top=5pt,bottom=5pt,
    before skip=0.9em,after skip=0.9em,
    fonttitle=\atl@headfont\bfseries,coltitle=#1!82!black,
    attach title to upper={\par\smallskip},
  },
}
\newcommand{\atl@newcallout}[3]{% env name, colour, string key
  \newtcolorbox{#1}[1][\atlstr{#3}]{atl callout=#2,title={##1}}}
\atl@newcallout{notebox}{atl-info}{note}
\atl@newcallout{warningbox}{atl-warning}{warning}
\atl@newcallout{tipbox}{atl-success}{tip}
\atl@newcallout{resultbox}{atl-result}{result}
\atl@newcallout{definitionbox}{atl-primary}{definition}
\atl@newcallout{examplebox}{atl-secondary-text}{example}

% ---- homework: problem / solution ---------------------------------------------
\newcounter{atlproblem}
\NewDocumentEnvironment{problem}{O{} O{}}{%
  \refstepcounter{atlproblem}%
  \begin{tcolorbox}[enhanced,breakable,sharp corners,boxrule=0.6pt,
    colframe=atl-primary,colback=atl-light,
    left=10pt,right=10pt,top=6pt,bottom=6pt,before skip=1.2em,after skip=0.6em,
    fonttitle=\atl@headfont\bfseries,coltitle=white,colbacktitle=atl-primary,
    toptitle=3pt,bottomtitle=3pt,
    title={\atlstr{problem}~\theatlproblem\ifstrempty{#1}{}{\,\textemdash\,#1}%
      \ifstrempty{#2}{}{\hfill\normalfont\small #2~\atlstr{points}}}]%
}{\end{tcolorbox}}
\NewDocumentEnvironment{solution}{O{}}{%
  \begin{tcolorbox}[enhanced,breakable,sharp corners,boxrule=0pt,frame hidden,
    colback=white,borderline \atl@start={2.4pt}{0pt}{atl-secondary},
    left=12pt,right=8pt,top=4pt,bottom=4pt,before skip=0.2em,after skip=1.2em,
    fonttitle=\atl@headfont\bfseries,coltitle=atl-secondary-text,
    attach title to upper={\par\smallskip},
    title={\atlstr{solution}\ifstrempty{#1}{}{\,\textemdash\,#1}}]%
}{\end{tcolorbox}}
\makeatother

% ============================================================================
% atl-code.tex -- syntax-highlighted code
%   key:   listing = darcula (default) | dark | bright | github
%   switch mid-document:  \atlsetlisting{bright}
%   plain listings:       \begin{lstlisting}[language=Python] ... \end{lstlisting}
%   framed with title:    \begin{codebox}{Python} ... \end{codebox}
%                         \begin{codebox}[title=Fast Fourier transform]{Python}
%   from a file:          \atlinputcode[title=main.c]{C}{../simulation/main.c}
% Languages: Python Matlab C C++ Verilog VHDL Assembler(x86masm) bash ... (any of listings').
% Code is always typeset left-to-right, also inside Persian documents.
% Legacy theme names dark_2026 and bright_2026 are accepted.
% ============================================================================
\makeatletter
% \atldeflisttheme{name}{bg=..,fg=..,kw=..,str=..,cmt=..,num=..,bi=..,deco=..,op=..,br=..}  (HTML hex)
\pgfkeys{/atl/lst/.cd,
  bg/.code={\definecolor{atl-\atl@curlt-bg}{HTML}{#1}},   fg/.code={\definecolor{atl-\atl@curlt-fg}{HTML}{#1}},
  kw/.code={\definecolor{atl-\atl@curlt-kw}{HTML}{#1}},   str/.code={\definecolor{atl-\atl@curlt-str}{HTML}{#1}},
  cmt/.code={\definecolor{atl-\atl@curlt-cmt}{HTML}{#1}}, num/.code={\definecolor{atl-\atl@curlt-num}{HTML}{#1}},
  bi/.code={\definecolor{atl-\atl@curlt-bi}{HTML}{#1}},   deco/.code={\definecolor{atl-\atl@curlt-deco}{HTML}{#1}},
  op/.code={\definecolor{atl-\atl@curlt-op}{HTML}{#1}},   br/.code={\definecolor{atl-\atl@curlt-br}{HTML}{#1}},
}
\newcommand{\atldeflisttheme}[2]{\def\atl@curlt{#1}\pgfkeys{/atl/lst/.cd,#2}}
\atldeflisttheme{darcula}{bg=2B2B2B,fg=A9B7C6,kw=CC7832,str=6A8759,cmt=808080,num=6897BB,bi=8888C6,deco=BBB529,op=A9B7C6,br=A9B7C6}
\atldeflisttheme{dark}{bg=141820,fg=C8D2DC,kw=FF8C32,str=78C878,cmt=646E78,num=64B4FF,bi=B496FF,deco=DCB432,op=C8D2DC,br=FF7878}
\atldeflisttheme{bright}{bg=FAFAFA,fg=282832,kw=B45014,str=148232,cmt=78828C,num=0064C8,bi=783CC8,deco=B47800,op=282832,br=C81E32}
\atldeflisttheme{github}{bg=0D1117,fg=C9D1D9,kw=FF7B72,str=A5D6FF,cmt=8B949E,num=79C0FF,bi=D2A8FF,deco=FFA657,op=C9D1D9,br=C9D1D9}

\newcommand{\atlsetlisting}[1]{%
  \def\atl@lt{#1}%
  \ifstrequal{#1}{dark_2026}{\def\atl@lt{dark}}{}%
  \ifstrequal{#1}{bright_2026}{\def\atl@lt{bright}}{}%
  \ifcsname\string\color@atl-\atl@lt-bg\endcsname\else
    \PackageError{atl}{Unknown listing theme `#1'}{Use darcula, dark, bright or github.}\fi
  \forcsvlist{\atl@lstrole}{bg,fg,kw,str,cmt,num,bi,deco,op,br}}
\newcommand{\atl@lstrole}[1]{\colorlet{atl-lst-#1}{atl-\atl@lt-#1}}

% numbers, operators and brackets get their own colour through literate rules
\lstdefinestyle{atl-lst}{
  basicstyle=\atl@ltr\ttfamily\footnotesize\color{atl-lst-fg},
  identifierstyle=\color{atl-lst-fg},
  keywordstyle=\color{atl-lst-kw},
  keywordstyle=[2]\color{atl-lst-bi},
  keywordstyle=[3]\color{atl-lst-deco},
  stringstyle=\color{atl-lst-str},
  commentstyle=\color{atl-lst-cmt}\itshape,
  numbers=left,numberstyle=\atl@ltr\tiny\color{atl-lst-cmt},numbersep=9pt,
  backgroundcolor=\color{atl-lst-bg},
  frame=single,rulecolor=\color{atl-lst-bg},framerule=0pt,framesep=7pt,
  xleftmargin=\dimexpr2.2em+7pt\relax,framexleftmargin=2.2em,xrightmargin=7pt,framexrightmargin=0pt,
  breaklines=true,breakatwhitespace=false,
  tabsize=4,showspaces=false,showstringspaces=false,upquote=true,
  columns=fullflexible,keepspaces=true,
  literate=
    *{0}{{\textcolor{atl-lst-num}{0}}}{1}
     {1}{{\textcolor{atl-lst-num}{1}}}{1}
     {2}{{\textcolor{atl-lst-num}{2}}}{1}
     {3}{{\textcolor{atl-lst-num}{3}}}{1}
     {4}{{\textcolor{atl-lst-num}{4}}}{1}
     {5}{{\textcolor{atl-lst-num}{5}}}{1}
     {6}{{\textcolor{atl-lst-num}{6}}}{1}
     {7}{{\textcolor{atl-lst-num}{7}}}{1}
     {8}{{\textcolor{atl-lst-num}{8}}}{1}
     {9}{{\textcolor{atl-lst-num}{9}}}{1}
     {=}{{\textcolor{atl-lst-op}{=}}}{1}
     {+}{{\textcolor{atl-lst-op}{+}}}{1}
     {-}{{\textcolor{atl-lst-op}{-}}}{1}
     {*}{{\textcolor{atl-lst-op}{*}}}{1}
     {/}{{\textcolor{atl-lst-op}{/}}}{1}
     {,}{{\textcolor{atl-lst-op}{,}}}{1}
     {:}{{\textcolor{atl-lst-op}{:}}}{1}
     {<}{{\textcolor{atl-lst-op}{<}}}{1}
     {>}{{\textcolor{atl-lst-op}{>}}}{1}
     {!}{{\textcolor{atl-lst-op}{!}}}{1}
     {?}{{\textcolor{atl-lst-op}{?}}}{1}
     {.}{{\textcolor{atl-lst-op}{.}}}{1}
     {;}{{\textcolor{atl-lst-op}{;}}}{1}
     {(}{{\textcolor{atl-lst-br}{(}}}{1}
     {)}{{\textcolor{atl-lst-br}{)}}}{1}
     {[}{{\textcolor{atl-lst-br}{[}}}{1}
     {]}{{\textcolor{atl-lst-br}{]}}}{1}
     {\{}{{\textcolor{atl-lst-br}{\{}}}{1}
     {\}}{{\textcolor{atl-lst-br}{\}}}}{1}
     {==}{{\textcolor{atl-lst-op}{==}}}{2}
     {!=}{{\textcolor{atl-lst-op}{!=}}}{2}
     {<=}{{\textcolor{atl-lst-op}{<=}}}{2}
     {>=}{{\textcolor{atl-lst-op}{>=}}}{2},
}

% a few extra built-ins per language (second keyword class)
\lstdefinelanguage{atlPython}[]{Python}{
  morekeywords=[2]{print,len,range,enumerate,open,zip,map,filter,sorted,sum,min,max,abs,int,float,str,list,dict,set,tuple,isinstance,super,type},
  morekeywords=[3]{self,cls,True,False,None}}
\lstdefinelanguage{atlC}[]{C}{
  morekeywords=[2]{printf,scanf,malloc,calloc,free,memcpy,memset,strlen,uint8_t,uint16_t,uint32_t,int8_t,int16_t,int32_t,size_t,bool},
  morekeywords=[3]{NULL,true,false}}
% friendly names -> our extended languages (plain lstlisting keeps the stock lexers)
\def\atl@l@Python{atlPython}\def\atl@l@python{atlPython}\def\atl@l@C{atlC}
\newcommand{\atl@maplang}[1]{\ifcsname atl@l@#1\endcsname\csname atl@l@#1\endcsname\else #1\fi}
\lstset{style=atl-lst}
\atldefault{listing}{darcula}
\edef\atl@tmp{\atlmeta{listing}}
\expandafter\atlsetlisting\expandafter{\atl@tmp}

% ---- framed code box ----------------------------------------------------------------
\newtcblisting{codebox}[2][]{%
  enhanced,breakable,listing only,
  listing options/.expanded={style=atl-lst,language=\atl@maplang{#2},numbers=none,frame=none,xleftmargin=0pt,framexleftmargin=0pt,backgroundcolor={}},
  colback=atl-lst-bg,colframe=atl-lst-bg!70!atl-lst-fg,boxrule=0.5pt,arc=3pt,
  borderline \atl@start={3pt}{0pt}{atl-accent},
  left=10pt,right=8pt,top=5pt,bottom=5pt,before skip=0.8em,after skip=0.8em,
  fonttitle=\atl@ltr\small\ttfamily,coltitle=atl-lst-fg,colbacktitle=atl-lst-bg!88!atl-lst-fg,
  title={#2},#1}
\newcommand{\atlinputcode}[3][]{%
  \tcbinputlisting{enhanced,breakable,listing only,
    listing file={#3},
    listing options/.expanded={style=atl-lst,language=\atl@maplang{#2},numbers=none,frame=none,xleftmargin=0pt,framexleftmargin=0pt,backgroundcolor={}},
    colback=atl-lst-bg,colframe=atl-lst-bg!70!atl-lst-fg,boxrule=0.5pt,arc=3pt,
    borderline \atl@start={3pt}{0pt}{atl-accent},
    left=10pt,right=8pt,top=5pt,bottom=5pt,before skip=0.8em,after skip=0.8em,
    fonttitle=\atl@ltr\small\ttfamily,coltitle=atl-lst-fg,colbacktitle=atl-lst-bg!88!atl-lst-fg,
    title={#2},#1}}
\makeatother

% ============================================================================
% atl-macros.tex -- text helpers and a small, general engineering/maths notation set
%
% Conventions (IEEE style): variables italic; descriptive subscripts, operators and
% units upright; vectors and matrices bold italic through ONE command each (\vect,
% \mat) so the look can be changed in one place. Units always go through siunitx:
%   \valunit{10}{\kilo\ohm}   \qty{2.2}{\micro\farad}   \num{1.5e-3}
% No LaTeX built-ins are redefined (\H \P \Re \Im \vec \div ...).
% Domain-specific packs (electronics, power, RF, ML ...) belong in preamble/local.tex.
% ============================================================================
\makeatletter
% ---- text helpers ----------------------------------------------------------------
\newcommand{\keyterm}[1]{\textbf{\textcolor{atl-primary}{#1}}}
\newcommand{\alert}[1]{\textbf{\textcolor{atl-warning}{#1}}}
\newcommand{\ok}[1]{\textbf{\textcolor{atl-success}{#1}}}
\newcommand{\highlight}[1]{\colorbox{atl-accent!28}{\strut #1}}
\newcommand{\muted}[1]{\textcolor{atl-muted}{#1}}
\newcommand{\todo}[1]{\textcolor{atl-warning}{\textbf{[TODO: #1]}}}

% ---- notation ---------------------------------------------------------------------
\newcommand{\qtag}[2]{\ensuremath{#1_{\mathrm{#2}}}}              % R_in
\newcommand{\qindex}[2]{\ensuremath{#1_{#2}}}                       % x_n
\newcommand{\qtagindex}[3]{\ensuremath{#1_{\mathrm{#2},#3}}}       % I_ref,k
\newcommand{\vect}[1]{\ensuremath{\bm{#1}}}
\newcommand{\mat}[1]{\ensuremath{\bm{#1}}}
\newcommand{\uvec}[1]{\ensuremath{\hat{\bm{#1}}}}
\newcommand{\phasor}[1]{\ensuremath{\tilde{#1}}}
\newcommand{\conj}[1]{\ensuremath{\overline{#1}}}
\newcommand{\trans}[1]{\ensuremath{{#1}^{\mathsf{T}}}}
\newcommand{\herm}[1]{\ensuremath{{#1}^{\mathsf{H}}}}
\newcommand{\pinv}[1]{\ensuremath{{#1}^{\dagger}}}
\DeclarePairedDelimiter{\abs}{\lvert}{\rvert}
\DeclarePairedDelimiter{\norm}{\lVert}{\rVert}
\DeclarePairedDelimiter{\inner}{\langle}{\rangle}
\newcommand{\ddiff}{\ensuremath{\mathop{}\!\mathrm{d}}}
\newcommand{\deriv}[2]{\ensuremath{\frac{\mathrm{d}#1}{\mathrm{d}#2}}}
\newcommand{\nderiv}[3]{\ensuremath{\frac{\mathrm{d}^{#3}#1}{\mathrm{d}#2^{#3}}}}
\newcommand{\pderiv}[2]{\ensuremath{\frac{\partial#1}{\partial#2}}}
\newcommand{\integral}[4]{\ensuremath{\int_{#1}^{#2}#3\,\mathrm{d}#4}}
\newcommand{\econst}{\ensuremath{\mathrm{e}}}
\newcommand{\junit}{\ensuremath{\mathrm{j}}}
\newcommand{\iunit}{\ensuremath{\mathrm{i}}}
\newcommand{\jw}{\ensuremath{\junit\omega}}
\newcommand{\N}{\ensuremath{\mathbb{N}}}
\newcommand{\Z}{\ensuremath{\mathbb{Z}}}
\newcommand{\Q}{\ensuremath{\mathbb{Q}}}
\newcommand{\R}{\ensuremath{\mathbb{R}}}
\newcommand{\C}{\ensuremath{\mathbb{C}}}
\DeclareMathOperator{\sgn}{sgn}
\DeclareMathOperator{\tr}{tr}
\DeclareMathOperator{\rank}{rank}
\DeclareMathOperator{\diag}{diag}
\DeclareMathOperator{\sinc}{sinc}
\DeclareMathOperator{\rect}{rect}
\DeclareMathOperator{\Expect}{\mathbb{E}}
\DeclareMathOperator{\Var}{Var}
\DeclareMathOperator{\Cov}{Cov}
\DeclareMathOperator{\Prob}{\mathbb{P}}
\DeclareMathOperator*{\argmin}{arg\,min}
\DeclareMathOperator*{\argmax}{arg\,max}

% ---- units (siunitx v3) -------------------------------------------------------------
\newcommand{\valunit}[3][]{\qty[#1]{#2}{#3}}
\DeclareSIUnit{\dBm}{dBm}
\DeclareSIUnit{\var}{var}
\DeclareSIUnit{\rpm}{rpm}
\DeclareSIUnit{\ppm}{ppm}
\DeclareSIUnit{\pu}{p.u.}
\makeatother

% ============================================================================
% atl-theorems.tex -- theorem-like environments (amsthm), coloured by the palette
%   theorem lemma proposition corollary (shared counter)  definition example remark  proof
% Labels come from the language layer; bodies are upright (Persian has no italic).
% ============================================================================
\makeatletter
\newtheoremstyle{atlplain}
  {0.8em}{0.8em}{\normalfont}{0pt}
  {\atl@headfont\bfseries\color{atl-primary}}{.}{0.5em}
  {\thmname{#1}\thmnumber{\,#2}\thmnote{\normalfont\ (#3)}}
\newtheoremstyle{atldef}
  {0.8em}{0.8em}{\normalfont}{0pt}
  {\atl@headfont\bfseries\color{atl-secondary-text}}{.}{0.5em}
  {\thmname{#1}\thmnumber{\,#2}\thmnote{\normalfont\ (#3)}}
\newtheoremstyle{atlremark}
  {0.6em}{0.6em}{\normalfont}{0pt}
  {\atl@headfont\bfseries\color{atl-muted}}{.}{0.5em}
  {\thmname{#1}\thmnumber{\,#2}\thmnote{\normalfont\ (#3)}}

\theoremstyle{atlplain}
\newtheorem{theorem}{\atlstr{theorem}}
\newtheorem{lemma}[theorem]{\atlstr{lemma}}
\newtheorem{proposition}[theorem]{\atlstr{proposition}}
\newtheorem{corollary}[theorem]{\atlstr{corollary}}
\theoremstyle{atldef}
\newtheorem{definition}{\atlstr{definition}}
\newtheorem{example}{\atlstr{example}}
\theoremstyle{atlremark}
\newtheorem{remark}{\atlstr{remark}}
\renewcommand{\proofname}{\atlstr{proof}}
\renewcommand{\qedsymbol}{\textcolor{atl-accent}{$\blacksquare$}}
\crefname{theorem}{\atlstr{theorem}}{\atlstr{theorem}}
\makeatother

% ============================================================================
% atl-titles.tex -- title pages.   Call  \atlmaketitle  right after \begin{document}.
%   key titlepage = block (default)  banner + info grid at the top of page 1
%                   cover             full cover page with palette artwork
%                   classic           centred journal-style title (essay, paper)
%                   none
%   keys used: title subtitle author affiliation studentid course coursecode instructor
%              supervisor advisor team partners experiment labsection term date
%              duedate submissiondate university department logo doctype category
%   infofields = comma list of keys shown in the info grid (labels come from the
%                language layer, so the same list works in English and Persian).
% Empty or unset keys are skipped.
% ============================================================================
\makeatletter
\atldefault{titlepage}{block}
\newcount\atl@cnt
\newlength{\atl@cellw}

\newcommand{\atl@doctype}{\atlifset{doctype}{\atlmeta{doctype}}{\atlstr{type-\atlmeta{category}}}}

\def\atl@amp{&}
\def\atl@endrow{\\[8pt]}
\newcommand{\atl@cellbody}[1]{%
  \global\advance\atl@cnt\@ne
  \parbox[t]{\atl@cellw}{\raggedright
    {\small\atl@headfont\color{atl-muted}\atlstr{#1}}\\[1pt]%
    {\bfseries\atlmeta{#1}}}%
  \ifodd\atl@cnt\expandafter\atl@amp\else\expandafter\atl@endrow\fi}
\newcommand{\atl@cell}[1]{\atlifset{#1}{\atl@cellbody{#1}}{}}
\newcommand{\atlinfogrid}[1]{%
  \begingroup
  \setlength{\atl@cellw}{0.465\linewidth}\global\atl@cnt\z@
  \edef\atl@l{#1}%
  \begin{tabular}{@{}p{\atl@cellw}@{\hspace{0.07\linewidth}}p{\atl@cellw}@{}}
    \expandafter\forcsvlist\expandafter\atl@cell\expandafter{\atl@l}%
    \ifodd\atl@cnt\mbox{}\fi
  \end{tabular}%
  \endgroup}

% ---- block: banner on top of page 1 -------------------------------------------------
\newcommand{\atl@title@block}{%
  \begingroup
  \begin{tcolorbox}[enhanced,sharp corners,frame hidden,colback=atl-primary,
      borderline south={4pt}{0pt}{atl-accent},
      left=15pt,right=15pt,top=13pt,bottom=13pt,before skip=0pt,after skip=0pt]
    {\atl@headfont\small\bfseries\color{white!70!atl-accent}\atl@doctype\par}\vspace{4pt}
    {\atl@headfont\bfseries\fontsize{24}{29}\selectfont\color{white}\atlmeta{title}\par}
    \atlifset{subtitle}{\vspace{4pt}{\atl@headfont\large\color{white!82!atl-primary}\atlmeta{subtitle}\par}}{}
  \end{tcolorbox}
  \vspace{1em}
  \atlinfogrid{\atlmeta{infofields}}\par
  \vspace{0.4em}{\color{atl-accent}\rule{\linewidth}{0.5pt}}\par\vspace{0.6em}
  \endgroup}

% ---- classic: journal-style --------------------------------------------------------
\newcommand{\atl@title@classic}{%
  \begin{center}
    {\atl@headfont\bfseries\LARGE\color{atl-primary}\atlmeta{title}\par}
    \atlifset{subtitle}{\vspace{4pt}{\large\color{atl-muted}\atlmeta{subtitle}\par}}{}
    \vspace{7pt}{\color{atl-accent}\rule{0.22\linewidth}{1.2pt}}\par\vspace{9pt}
    \atlifset{author}{{\large\atlmeta{author}\par}}{}
    \atlifset{affiliation}{\vspace{2pt}{\small\color{atl-muted}\atlmeta{affiliation}\par}}{}
    \atlifset{date}{\vspace{2pt}{\small\atlmeta{date}\par}}{}
  \end{center}\vspace{0.4em}}

% ---- cover: full page ----------------------------------------------------------------
\newcommand{\atl@cover@art}{%
  \begin{tikzpicture}[remember picture,overlay]
    \fill[atl-tint] (current page.south west) rectangle (current page.north east);
    % top ribbons (back to front)
    \path[fill=atl-accent!85!atl-tint] (current page.north west)
      -- ([yshift=-6.6cm]current page.north west)
      to[out=-12,in=188] ([yshift=-2.9cm]current page.north east)
      -- (current page.north east) -- cycle;
    \path[fill=atl-secondary!90!atl-tint] (current page.north west)
      -- ([yshift=-5.7cm]current page.north west)
      to[out=-16,in=190] ([yshift=-2.1cm]current page.north east)
      -- (current page.north east) -- cycle;
    \path[fill=atl-primary] (current page.north west)
      -- ([yshift=-4.7cm]current page.north west)
      to[out=-20,in=192] ([yshift=-1.3cm]current page.north east)
      -- (current page.north east) -- cycle;
    \path[fill=white,fill opacity=0.10] (current page.north west)
      -- ([yshift=-3.1cm]current page.north west)
      to[out=-24,in=195] ([yshift=-0.5cm]current page.north east)
      -- (current page.north east) -- cycle;
    % bottom ribbons
    \path[fill=atl-accent!85!atl-tint] (current page.south east)
      -- ([yshift=2.9cm]current page.south east)
      to[out=172,in=-8] ([yshift=1.2cm]current page.south west)
      -- (current page.south west) -- cycle;
    \path[fill=atl-secondary!90!atl-tint] (current page.south east)
      -- ([yshift=2.2cm]current page.south east)
      to[out=170,in=-10] ([yshift=0.8cm]current page.south west)
      -- (current page.south west) -- cycle;
    \path[fill=atl-primary] (current page.south east)
      -- ([yshift=1.5cm]current page.south east)
      to[out=168,in=-12] ([yshift=0.4cm]current page.south west)
      -- (current page.south west) -- cycle;
  \end{tikzpicture}}

\newcommand{\atl@title@cover}{%
  \begin{titlepage}
    \thispagestyle{empty}\atl@cover@art
    \vspace*{3.4cm}
    \begin{center}
      \atlifset{logo}{\includegraphics[height=2.8cm]{\atlmeta{logo}}\par\vspace{0.7cm}}{}
      \atlifset{university}{{\atl@headfont\Large\bfseries\color{atl-primary}\atlmeta{university}\par}}{}
      \atlifset{department}{\vspace{2pt}{\atl@headfont\large\color{atl-secondary-text}\atlmeta{department}\par}}{}
      \vspace{1.6cm}
      {\atl@headfont\small\bfseries\color{atl-accent-text}\atl@doctype\par}\vspace{8pt}
      {\atl@headfont\bfseries\fontsize{30}{36}\selectfont\color{atl-primary}\atlmeta{title}\par}
      \atlifset{subtitle}{\vspace{8pt}{\atl@headfont\Large\color{atl-muted}\atlmeta{subtitle}\par}}{}
      \vspace{14pt}{\color{atl-accent}\rule{0.28\linewidth}{2pt}}\par
    \end{center}
    \vfill
    \atlinfogrid{\atlmeta{infofields}}
    \vspace{2.6cm}
  \end{titlepage}}

\newcommand{\atlmaketitle}{%
  \atlifeq{titlepage}{cover}{\atl@title@cover}{%
  \atlifeq{titlepage}{classic}{\atl@title@classic}{%
  \atlifeq{titlepage}{none}{}{\atl@title@block}}}}

\AtBeginDocument{%
  \hypersetup{unicode=true,pdftitle={\atlmeta{title}},pdfauthor={\atlmeta{author}}}}
\makeatother

% ============================================================================
% atl-bib.tex -- bibliography (BibTeX + IEEEtran, numeric [1], [2], ...)
%   \atlbibliography{refs}      refs.bib in the project; heading and ToC entry included
%   key bibstyle = IEEEtran (default)
% The list is typeset left-to-right (also in Persian documents) because entries are
% Latin. Cite with \cite{key}. latexmk runs BibTeX automatically.
% ============================================================================
\makeatletter
\atldefault{bibstyle}{IEEEtran}
\renewenvironment{thebibliography}[1]{%
  \atlunnumbered{\atlstr{references}}%
  \begingroup\atl@ltr\small
  \list{\@biblabel{\@arabic\c@enumiv}}{%
    \settowidth\labelwidth{\@biblabel{#1}}%
    \leftmargin\labelwidth \advance\leftmargin\labelsep
    \usecounter{enumiv}\let\p@enumiv\@empty
    \renewcommand\theenumiv{\@arabic\c@enumiv}%
    \itemsep=0.2em}%
  \sloppy\clubpenalty4000\widowpenalty4000\sfcode`\.\@m}
 {\def\@noitemerr{\@latex@warning{Empty `thebibliography' environment}}%
  \endlist\endgroup}
\newcommand{\atlbibliography}[1]{%
  \edef\atl@bs{\atlmeta{bibstyle}}%
  \expandafter\bibliographystyle\expandafter{\atl@bs}%
  \bibliography{#1}}
\makeatother


\makeatletter % (each module ends with \makeatother)
% names LaTeX prints itself (Contents, List of Figures, ...) come from the language layer
\AtBeginDocument{%
  \renewcommand{\contentsname}{\atlstr{contents}}%
  \renewcommand{\listfigurename}{\atlstr{listoffigures}}%
  \renewcommand{\listtablename}{\atlstr{listoftables}}%
  \renewcommand{\figurename}{\atlstr{figure}}%
  \renewcommand{\tablename}{\atlstr{table}}%
  \renewcommand{\appendixname}{\atlstr{appendix}}%
  \@ifundefined{refname}{\newcommand{\refname}{\atlstr{references}}}{\renewcommand{\refname}{\atlstr{references}}}%
  \@ifundefined{bibname}{}{\renewcommand{\bibname}{\atlstr{references}}}%
  \@ifundefined{abstractname}{}{\renewcommand{\abstractname}{\atlstr{abstract}}}}
\hypersetup{linktoc=page}
\makeatother
\makeatother
                           % ۴. موتور قالب
% ماکروها و تنظیمات شخصی شما؛ پس از موتور قالب بارگذاری می‌شود.
% \atlsetpalette{green}                 % تغییر پالت رنگ از همین‌جا
% \atlsetlisting{bright}                % تغییر پوستهٔ کد (darcula, dark, bright, github)
% \newcommand{\A}{\mat{A}}

% ---- ماکروهای این تمرین ----
\definecolor{navy}{HTML}{1F3A68}
\definecolor{accent}{HTML}{B3401E}
\definecolor{lane1}{HTML}{E8F0FE}
\definecolor{lane2}{HTML}{FDF1E3}
\newcommand{\code}[1]{\lr{\texttt{#1}}}                      % نام فراخوانی، ثبات، مسیر
\newcommand{\src}[1]{\par\begingroup\color{gray}\small\noindent\textbf{منبع:}~#1\par\endgroup}
                    % ۵. ماکروها و تنظیمات شخصی شما

\begin{document}
\atlmaketitle

\begin{notebox}[راهنمای خواندن]
در انتهای هر بخش «منبع» آمده است. عبارت «اسلاید» یعنی پاسخ بر پایهٔ اسلایدهای فصل ۱ (مقدمه) و فصل ۲ (ساختار سیستم‌عامل) نوشته شده است و عبارت «تکمیلی» یعنی آن مطلب در اسلایدها نبوده و پاسخ از دانش درسی (کتاب سیلبرشاتس، مستندات لینوکس و مقالات مربوط) تکمیل شده است.
\end{notebox}

\section{مجازی‌سازی}

\subsection*{بخش آ – اجرای مستقیم چند سیستم‌عامل بدون لایهٔ مجازی‌سازی}

هر سیستم‌عاملِ مهمان برای کار کردن فرض می‌کند که تنها مالک ماشین است و در حالت هسته اجرا می‌شود. سخت‌افزار در زمان بوت در حالت هسته شروع می‌کند و فقط در این حالت می‌توان دستورهای ممتاز (مثل تنظیم زمان‌سنج، مدیریت وقفه، کنترل \lr{I/O} و پاک‌کردن حافظه) را اجرا کرد. بنابراین یک سیستم‌عامل مهمان از سخت‌افزار این انتظارها را دارد:
\begin{itemize}
  \item \textbf{پردازنده:} اجرای آزاد دستورهای ممتاز، تنظیم زمان‌سنج و کنترل‌گر وقفه، ماسک‌کردن وقفه‌ها، تعویض جدول وقفه (\lr{IDT}) و داشتن \lr{CPU} به‌صورت انحصاری. زمان‌بندی را خودش انجام می‌دهد و هیچ دلیلی ندارد پردازنده را به سیستم‌عامل دیگری بدهد.
  \item \textbf{حافظه:} کل حافظهٔ فیزیکی (از آدرس صفر و پیوسته) را از آنِ خود می‌داند، جدول صفحه و \lr{MMU} را خودش مقداردهی می‌کند و می‌تواند هر آدرس فیزیکی را بخواند یا بنویسد.
  \item \textbf{ورودی/خروجی:} دسترسی مستقیم به ثبات‌های کنترل‌گر دستگاه‌ها، درگاه‌ها، خطوط وقفه و \lr{DMA}؛ و انتظار دارد دستگاه‌ها دقیقاً در همان حالتی باشند که خودش رها کرده است.
\end{itemize}

اگر چند سیستم‌عامل هم‌زمان چنین کنترلی داشته باشند، مشکلات زیر پیش می‌آید:
\begin{itemize}
  \item \textbf{تداخل حافظه:} هر هسته روی حافظهٔ دیگری می‌نویسد؛ یک سیستم‌عامل می‌تواند هستهٔ سیستم‌عامل دیگر را پاک کند (همان سناریوی «برنامهٔ کاربر می‌تواند سیستم‌عامل را پاک کند»، این بار میان سیستم‌عامل‌ها).
  \item \textbf{تداخل در \lr{I/O}:} چند سیستم‌عامل هم‌زمان در یک دستگاه می‌نویسند و دستورهای متناقض به یک کنترل‌گر می‌رسد؛ معلوم هم نیست یک وقفه باید به کدام سیستم‌عامل تحویل شود.
  \item \textbf{تداخل در زمان‌بندی:} زمان‌سنج و کنترل‌گر وقفه یکی است و هر سیستم‌عامل آن را به شکل خودش تنظیم می‌کند؛ تنظیمِ دومی تنظیمِ اولی را خراب می‌کند. هیچ‌کس نمی‌تواند \lr{CPU} را از سیستم‌عاملی که وقفه‌ها را غیرفعال کرده یا در حلقهٔ بی‌نهایت است پس بگیرد؛ تقسیم عادلانهٔ منابع ناممکن است.
  \item \textbf{اجرای نادرست:} همه انتظار دارند توسط میان‌افزار بوت شوند و جدول وقفه، جدول صفحه و ثبات‌های کنترلی را در اختیار بگیرند؛ در هر لحظه فقط یک مجموعه از این وضعیت‌ها روی سخت‌افزار می‌تواند فعال باشد و تعویض میان سیستم‌عامل‌ها (ذخیره و بازیابی کل وضعیت ماشین) کار هیچ‌کدام نیست.
  \item \textbf{امنیت و پایداری:} هیچ مرز حفاظتی وجود ندارد؛ خطا یا بدافزار در یک سیستم‌عامل کل ماشین را از کار می‌اندازد یا اطلاعات بقیه را می‌خواند.
\end{itemize}

\textbf{چرا لایهٔ واسط لازم است؟} چون باید یک مرجع واحد وجود داشته باشد که کنترل واقعی منابع را در دست دارد و بقیه را به رعایت قواعد مجبور می‌کند. این لایه همان ناظر ماشین مجازی (\lr{Hypervisor/VMM}) است؛ نرم‌افزار (یا سخت‌افزار)ی است که ماشین‌های مجازی را می‌سازد و اجرا می‌کند (نوع ۱ مستقیماً روی سخت‌افزار) و برای هر مهمان این توهم را می‌سازد که کامپیوتر خصوصی خودش را دارد. وظایف جداسازی آن:
\begin{itemize}
  \item \textbf{پردازنده:} زمان‌بندی \lr{vCPU}ها روی هسته‌های واقعی با زمان‌سنجی که فقط ناظر کنترل می‌کند (پیش‌دستی اجباری؛ به همان دلیل که سیستم‌عامل با زمان‌سنج از انحصار پردازه‌ها جلوگیری می‌کند).
  \item \textbf{حافظه:} نگاشت «فیزیکی مهمان $\rightarrow$ فیزیکی واقعی» با جدول صفحهٔ سایه یا جدول صفحهٔ تودرتو، تخصیص و سهمیه‌بندی حافظه و جلوگیری از دسترسی به حافظهٔ دیگران.
  \item \textbf{ورودی/خروجی:} ارائهٔ دستگاه‌های مجازی (شبیه‌سازی‌شده یا نیمه‌مجازی) یا تخصیص اختصاصی دستگاه همراه با \lr{IOMMU}، و تحویل وقفه‌های مجازی به مهمان درست.
  \item \textbf{حفاظت از خود و مهار خطا:} دستورهای ممتاز فقط در دست ناظر است؛ خرابی یا آلودگی یک مهمان فقط خودش را از کار می‌اندازد و میزبان و مهمان‌های دیگر محافظت می‌شوند.
\end{itemize}

\subsection*{بخش ب – اجرای هستهٔ مهمان مستقیماً در حالت هستهٔ واقعی}

اگر ناظر اجازه دهد هستهٔ مهمان در حالت هستهٔ واقعی اجرا شود، عملاً ناظر کنترل ماشین را از دست داده است:
\begin{itemize}
  \item \textbf{امنیت و دسترسی غیرمجاز:} مهمان می‌تواند هر دستور ممتازی را اجرا کند؛ پس می‌تواند هر آدرس فیزیکی (حافظهٔ ناظر و مهمان‌های دیگر) را بخواند و بنویسد، جدول صفحه و \lr{CR3} را عوض کند یا با \lr{DMA} حافظهٔ دیگران را بخواند. هیچ سازوکاری نمی‌تواند دسترسی غیرمجاز را به تله بیندازد.
  \item \textbf{جداسازی و اثرگذاری روی دیگران:} مهمان می‌تواند جدول وقفه یا کد ناظر را بازنویسی کند و حتی ناظر را دور بزند. چون وقفه‌ها را ماسک و زمان‌سنج را دوباره‌تنظیم می‌کند، ناظر هرگز کنترل \lr{CPU} را پس نمی‌گیرد؛ مهمانِ دارای حلقهٔ بی‌نهایت بقیه را گرسنه نگه می‌دارد (همان مسئلهٔ «\lr{process hogging}» که زمان‌سنج برای جلوگیری از آن وجود دارد).
  \item \textbf{کنترل و تخصیص منابع:} ناظر نمی‌تواند سهمیه بگذارد یا محدودیتی اعمال کند، چون مهمان خود را مالک تمام حافظه و تمام \lr{CPU}ها می‌داند. دستگاه‌های \lr{I/O} به‌طور مستقیم و هم‌زمان توسط چند مهمان برنامه‌ریزی می‌شوند و وقفه‌ها و پیکربندی‌ها با هم تداخل می‌کنند؛ اشتراک‌گذاری کنترل‌شدهٔ دستگاه ناممکن است.
  \item \textbf{پایداری کل سیستم:} خطا در هستهٔ مهمان (اشارهگر خراب، دستور \lr{HLT}، خطای سه‌گانه، پیکربندی غلط سخت‌افزار) مستقیماً ماشین واقعی را متوقف یا بازنشانی می‌کند و فقط به همان مهمان محدود نمی‌ماند. ناظر هم نمی‌تواند وضعیت مهمان را ببیند یا ذخیره کند، پس فعالیت‌هایی مثل عکس‌گیری (\lr{snapshot}) و انتقال زندهٔ ماشین مجازی هم ممکن نیست.
\end{itemize}

نتیجه: هستهٔ مهمان باید از حالت هستهٔ واقعی «تنزل» پیدا کند (مثلاً به حالت کاربر، مانند بخش بعد) یا سخت‌افزار حالت جداگانه‌ای برای مهمان فراهم کند (بخش ه).

\subsection*{بخش ج – روش «تله و شبیه‌سازی» (\lr{Trap-and-Emulate})}

\textbf{روش کار.} ناظر در حالت هستهٔ واقعی اجرا می‌شود و هستهٔ مهمان را در حالت کاربر واقعی اجرا می‌کند؛ مهمان فقط یک «حالت کاربر مجازی» و «حالت هستهٔ مجازی» دارد که هر دو روی حالت کاربر واقعی اجرا می‌شوند. مراحل:
\begin{itemize}
  \item دستورهای معمولی مهمان، چه کد کاربر و چه کد هستهٔ مهمان، مستقیماً و با سرعت تقریباً کامل روی \lr{CPU} اجرا می‌شوند.
  \item وقتی مهمان دستور ممتازی اجرا می‌کند (مثلاً تغییر جدول صفحه، غیرفعال‌کردن وقفه، دستور \lr{I/O})، چون \lr{CPU} در حالت کاربر است، استثنا (تله) رخ می‌دهد و کنترل با ذخیرهٔ وضعیت به ناظر منتقل می‌شود.
  \item ناظر دستور را رمزگشایی می‌کند و اثر آن را روی وضعیت مجازی مهمان (ثبات‌ها و بیت حالت مجازی، جدول صفحهٔ مجازی، دستگاه مجازی) شبیه‌سازی می‌کند و نه روی سخت‌افزار واقعی.
  \item ناظر وضعیت مهمان را بازیابی می‌کند و اجرای مهمان را از دستور بعدی ادامه می‌دهد.
\end{itemize}

\textbf{محدودیت‌ها در معماری‌هایی که برای مجازی‌سازی طراحی نشده‌اند.}
\begin{itemize}
  \item \textbf{دستورهای حساسِ غیرممتاز:} در برخی معماری‌ها (مثل \lr{x86} پیش از افزوده‌شدن پشتیبانی سخت‌افزاری) دستورهایی هستند که وضعیت یا پیکربندی سیستم را می‌خوانند یا تغییر می‌دهند ولی ممتاز نیستند؛ مثل \code{POPF} و \code{SGDT}. این دستورها در حالت کاربر تله ایجاد نمی‌کنند، پس ناظر هرگز مطلع نمی‌شود و نتیجهٔ نادرست به مهمان می‌رسد (تفصیل در بخش د).
  \item مجبور می‌شویم از راه‌حل‌های سنگین‌تر مانند بازنویسی دودویی (ترجمهٔ کد مهمان) یا نیمه‌مجازی‌سازی (تغییر دادن هستهٔ مهمان) استفاده کنیم.
  \item \textbf{هزینهٔ زیاد تله‌ها:} هستهٔ مهمان تعداد زیادی دستور ممتاز اجرا می‌کند و هر کدام به تله می‌افتد؛ این هم باعث افت کارایی می‌شود.
\end{itemize}

\textbf{چرا انتقال‌های زیاد کارایی را کم می‌کند؟} هر انتقال شامل ایجاد استثنا، ذخیرهٔ وضعیت مهمان، سوئیچ به زمینهٔ ناظر، رمزگشایی و شبیه‌سازی دستور، بازیابی وضعیت و بازگشت است. اگر دستور ممتاز روی ماشین واقعی حدود یک چرخهٔ ساعت طول می‌کشد، همین دستور اینجا صدها تا هزاران چرخه هزینه دارد. افزون بر آن، هر انتقال خط لولهٔ \lr{CPU} را خالی می‌کند و کش و \lr{TLB} را آلوده می‌کند. چون هستهٔ مهمان مدام جدول صفحه را به‌روز می‌کند، وقفه‌ها را فعال و غیرفعال می‌کند و به ثبات‌های دستگاه دسترسی می‌دهد (و هر عملیات \lr{I/O} چندین دسترسی به ثبات دارد)، این هزینه چندین برابر می‌شود.

\subsection*{بخش د – ناتوانی معماری در تله‌انداختن همهٔ رفتارهای حساس}

بر اساس شرط پوپک–گلدبرگ، معماری‌ای برای مجازی‌سازی با تله و شبیه‌سازی مناسب است که مجموعهٔ دستورهای حساس (آن‌هایی که وضعیت و پیکربندی منابع را می‌خوانند یا عوض می‌کنند یا بسته به حالت رفتار متفاوت دارند) زیرمجموعهٔ دستورهای ممتاز باشد. اگر چنین نباشد:
\begin{itemize}
  \item \textbf{ایجاد مشکل در پیاده‌سازی:} مهمان دستور حساس را در حالت کاربر مستقیماً اجرا می‌کند و ناظر خبردار نمی‌شود. یا دستور بی‌صدا نتیجهٔ اشتباه می‌دهد (مثلاً \code{POPF} در حالت کاربر تغییر پرچم وقفه را نادیده می‌گیرد، در حالی که مهمان فکر می‌کند وقفه‌ها را غیرفعال کرده است و منطق هسته‌اش خراب می‌شود)، یا اطلاعات واقعی ماشین را به مهمان نشان می‌دهد (مثلاً \code{SGDT} آدرس واقعی جدول توصیف‌گر سراسری را برمی‌گرداند و نه آدرس مجازی). در نتیجه ناظر دیگر نمی‌تواند به مهمان توهم «یک ماشین کامل» بدهد (ویژگی هم‌ارزی).
  \item \textbf{نقض جداسازی:} اگر اثر دستور روی حالت واقعی سخت‌افزار برود، یک مهمان بر وضعیت مشترک ماشین اثر می‌گذارد و بقیهٔ ماشین‌های مجازی و حتی ناظر را تحت تأثیر قرار می‌دهد. اگر دستور اطلاعات بخواند، اطلاعات ناظر یا مهمان‌های دیگر نشت می‌کند و مهمان می‌فهمد روی ماشین مجازی اجرا می‌شود. در هر دو حالت، ناظر کنترل کامل منابع را ندارد (ویژگی کنترل منابع) و مرز حفاظتی شکسته می‌شود.
  \item \textbf{چرا تشخیص صحیح عملیات حساس نیازمندی اصلی است؟} چون هر سه ویژگی مورد نیاز مجازی‌سازی، یعنی وفاداری (مهمان همان رفتار ماشین واقعی را ببیند)، کنترل کامل منابع توسط ناظر و کارایی (بیشتر دستورها مستقیم اجرا شوند)، وقتی برقرار است که ناظر دقیقاً در لحظه‌های لازم دخالت کند. تشخیص ناقص یعنی یا ناظر دخالت نمی‌کند (خطای رفتار و نقض جداسازی)، یا مجبور است همه‌چیز را شبیه‌سازی کند (کارایی ضعیف).
\end{itemize}

\subsection*{بخش ه – پشتیبانی سخت‌افزاری از مجازی‌سازی}

پردازنده‌های جدید (مثل \lr{Intel VT-x} و \lr{AMD-V}) یک حالت اجرای جدید اضافه می‌کنند که «حالت میزبان/ناظر» در برابر «حالت مهمان» نامیده می‌شود. ناظر در حالت میزبان اجرا می‌شود و مهمان در حالت مهمان؛ هستهٔ مهمان در حالت مهمان، حالت هستهٔ مجازی خودش را دارد و می‌تواند بیشتر دستورهای ممتاز را مستقیماً اجرا کند. هر رویدادی که ناظر باید در آن دخالت کند (مثلاً دسترسی به دستگاه، تغییر بعضی ثبات‌های کنترلی، خطای صفحهٔ تودرتو) با تنظیمات ناظر باعث «خروج از ماشین مجازی» (\lr{VM exit}) می‌شود و \lr{CPU} به‌طور خودکار وضعیت مهمان را در ساختار سخت‌افزاری (\lr{VMCS/VMCB}) ذخیره و وضعیت ناظر را بارگذاری می‌کند.

\textbf{تأثیر بر طراحی و عملکرد ناظر:}
\begin{itemize}
  \item \textbf{ساده‌تر شدن طراحی:} ناظر دیگر لازم نیست کد مهمان را اسکن یا بازنویسی کند و کد هسته‌ی مهمان را هم تغییر نمی‌دهد (بدون بازنویسی دودویی و بدون نیمه‌مجازی‌سازی). ناظر به یک حلقهٔ رویدادمحور تبدیل می‌شود: ورود به مهمان، خروج با علت مشخص، پردازش علت (شبیه‌سازی دستگاه یا مدیریت منابع) و ورود دوباره.
  \item چون همهٔ رفتارهای حساس دیگر قابل‌تله‌انداختن‌اند، مشکل دستورهای حساسِ غیرممتاز از بین می‌رود.
  \item ذخیره و بازیابی وضعیت مهمان را سخت‌افزار انجام می‌دهد و ناظر مجبور نیست آن را دستی ساخته و نگه دارد.
  \item جدول صفحهٔ تودرتو (\lr{EPT/NPT}) ترجمهٔ آدرس مهمان را به سخت‌افزار می‌سپارد و ناظر از نگهداری جدول صفحهٔ سایه خلاص می‌شود.
  \item \textbf{کارایی:} مهمان‌های بدون تغییر با سرعت نزدیک به بومی اجرا می‌شوند و فقط رویدادهای واقعاً حساس هزینهٔ خروج را می‌پردازند.
\end{itemize}


\section{وقفه به‌عنوان ابزار حمله}

\textbf{بخش اول: حملهٔ \lr{Heckler}}

\subsection*{بخش آ – چرا گردانندهٔ فراخوانی سیستمی در اجرای ناخواسته به‌صورت وقفه خطرناک است؟}

تله و وقفه از یک جدول بردار و یک مسیر مشترک ورود به هسته استفاده می‌کنند. گردانندهٔ فراخوانی سیستمی بر چند فرض بنا شده است:
\begin{itemize}
  \item کسی با اجرای عمدی دستور تله (مثل \code{int 0x80}) آن را صدا زده است و نشانی بازگشت، حالت و ثبات‌ها بر اساس قرارداد فراخوانی آماده شده‌اند.
  \item هسته «دستور وقفه‌دهنده» را بررسی می‌کند تا نوع را بفهمد و سپس پارامترهای ثبات را کنترل و فراخوانی را اجرا می‌کند؛ یعنی منبع و نیت فراخوانی را مفروض می‌گیرد.
  \item فراخواننده یک پردازهٔ کاربر است که خودش از اثرهای فراخوانی (فقط نتیجه در یک ثبات) آگاه است.
\end{itemize}

وقتی همین گردانندهٔ همین بردار با یک وقفهٔ تزریقی، یعنی ناهم‌زمان و در یک لحظهٔ دلخواه، اجرا شود، هیچ‌کدام از این فرض‌ها برقرار نیست. سخت‌افزار برای وقفه‌های خارجی (و تزریق‌شده) بررسی سطح دسترسی دروازه را انجام نمی‌دهد (این بررسی فقط برای دستور نرم‌افزاری \code{int n} است)، و مسیر مشترک هم نمی‌تواند تشخیص دهد که وقفه از یک دستور واقعی آمده است یا از بیرون. در نتیجه هسته محتوای تصادفی ثبات‌های پردازهٔ قربانی را به‌عنوان شمارهٔ فراخوانی و آرگومان می‌پذیرد، یک فراخوانی سیستمی دلخواه را با اختیارات پردازه اجرا می‌کند و نتیجه را روی ثبات‌هایی می‌نویسد که پردازه در حال استفاده از آن‌هاست. مهاجم زمان تزریق را طوری انتخاب می‌کند که اجرای پردازه را به نفع خود تغییر دهد.

\subsection*{بخش ب – اهمیت مقدار ثبات‌ها در لحظهٔ حمله}

هسته شمارهٔ فراخوانی را از ثبات \code{rax} و آرگومان‌ها را از ثبات‌های \code{rdi}، \code{rsi}، \code{rdx}، \code{r10}، \code{r8} و \code{r9} می‌خواند و مقدار برگشتی را هم در \code{rax} می‌گذارد. نتیجه:
\begin{itemize}
  \item مقدار ثبات‌ها در لحظهٔ تزریق از نظر مهاجم «ورودیِ فراخوانی» است. مهاجم با انتخاب لحظه می‌تواند شمارهٔ فراخوانی و آرگومان‌های مطلوب را بسازد، چون ثبات‌ها محتوای محاسبه‌ای را دارند که پردازه در همان لحظه انجام می‌دهد.
  \item نتیجهٔ فراخوانی ساختگی روی \code{rax} می‌نشیند. اگر \code{rax} در آن لحظه مقدار برگشتی تابعی مثل بررسی گذرواژه باشد (تابع در \code{rax} نتیجه می‌گذارد)، مقدار برگشتی با نتیجهٔ فراخوانی ساختگی (مثلاً صفر یعنی موفقیت) جایگزین می‌شود و شاخهٔ برنامه عوض می‌شود. روش دور زدن احراز هویت در \lr{sudo} و \lr{OpenSSH} همین است.
  \item چون سیستم‌عامل مهمان محرمانه است، ناظر نمی‌تواند حافظه و ثبات‌ها را ببیند. او حالت ثبات‌ها را با زمان‌بندی (مثلاً شمارش وقفه‌ها یا رویدادها پیش از ورود به تابع حساس) حدس می‌زند و بنابراین حمله به زمان‌بندی دقیق و مقدار ثبات‌ها وابسته است.
\end{itemize}

\subsection*{بخش ج – ویژگی روال رسیدگی امن و یک راه دفاع}

روالی در هر لحظه بی‌خطر است که (۱) وضعیت‌ ناشی از اجرای خودش را به‌طور کامل ذخیره و بازیابی کند و هیچ اثری روی برنامهٔ وقفه‌خورده نگذارد؛ (۲) فرضی دربارهٔ محتوای ثبات‌ها یا زمان فراخوانی نداشته باشد و ورودی‌هایش را اعتبارسنجی کند؛ (۳) تنها از مسیر و منبع مجاز قابل ورود باشد؛ یعنی بداند به‌دلیل یک دستور واقعی در حالت کاربر فراخوانده شده است؛ (۴) بازورود‌پذیر باشد و با ماسک‌کردن مناسب وقفه‌ها به وضعیت ناسازگار نرسد.

\textbf{راه دفاع پیشنهادی:} چند لایه ممکن است:
\begin{itemize}
  \item \textbf{سطح هسته:} اگر مسیر \code{int 0x80} (فراخوانی سیستمی قدیمی ۳۲‌بیتی) لازم نیست، غیرفعال شود (برای مثال با تنظیمات \code{ia32\_emulation}) تا بردار \code{0x80} گردانندهٔ فعال نداشته باشد. لینوکس برای مهمان‌های \lr{Intel TDX} و \lr{AMD SEV} این قابلیت را به‌صورت پیش‌فرض غیرفعال می‌کند.
  \item \textbf{سطح سخت‌افزار و ناظر:} محدودکردن بردارهایی که ناظر اجازه دارد تزریق کند (مثل ویژگی تزریق محدود/جایگزین در \lr{AMD SEV-SNP} و بررسی تزریق در \lr{Intel TDX})، تا بردارهای ویژه‌ای مثل \code{0x80} و بردارهای استثنا قابل‌تزریق نباشند.
  \item \textbf{سطح گرداننده:} اجرای گردانندهٔ فراخوانی سیستمی فقط وقتی مجاز باشد که با دستور واقعی در حالت کاربر آمده باشد (برای نمونه ثبت منشأ ورود در مسیر مشترک و رد وقفه‌های بدون منشأ معتبر).
\end{itemize}

\textbf{بخش دوم: نشت اطلاعات از طریق وقفه‌ها (حملهٔ شمارنده در مرورگر)}

\subsection*{بخش د – زبانه‌ها چگونه از هم باخبر می‌شوند و چرا هر وب‌گاه الگوی خودش را دارد؟}

دو زبانه حافظهٔ مشترک ندارند، ولی \textbf{سخت‌افزار و هسته را به اشتراک گذاشته‌اند}. زبانهٔ مهاجم حلقه‌ای می‌چرخاند که یک شمارنده را زیاد می‌کند و در هر بازهٔ زمانی کوتاه (مثلاً چند میلی‌ثانیه) مقدار آن را ثبت می‌کند. هر وقت \lr{CPU} (یا همان هستهٔ پردازنده‌ای که حلقه روی آن است) به کار دیگری مشغول شود، شمارنده کمتر زیاد می‌شود و در دنباله نقطهٔ افتادگی می‌ماند. دو عامل اصلی این افتادگی‌ها هستند:
\begin{itemize}
  \item \textbf{وقفه‌ها:} هر وقفه (کارت شبکه، پردازندهٔ گرافیکی، دیسک، زمان‌سنج، وقفهٔ بین‌هسته‌ای) \lr{CPU} را متوقف می‌کند تا گرداننده اجرا شود؛ پردازنده با هر وقفه وضعیت را ذخیره می‌کند و به گردانندهٔ مربوط می‌رود. مدت و زمان این وقفه‌ها از رفتار «زبانهٔ قربانی» می‌آید.
  \item \textbf{رقابت بر منابع مشترک:} کش، پهنای باند حافظه، منابع داخلی هسته و فرکانس پردازنده.
\end{itemize}

هر وب‌گاه الگوی خاص خودش را می‌سازد، چون اندازه، تعداد و ترتیب منابع و درخواست‌های شبکه، زمان‌بندی دریافت بسته‌ها، رندر و اجرای جاوااسکریپت (و در نتیجه فعالیت \lr{GPU} و \lr{CPU}) برای هر وب‌گاه متفاوت است. نتیجه یک «اثر انگشت» زمانی از وقفه‌ها و بار \lr{CPU} است که شمارنده آن را به‌صورت یک سری زمانی ثبت می‌کند و مدل یادگیری ماشین (مثل شبکهٔ عصبی کانولوشنی) آن را به نام وب‌گاه نگاشت می‌کند.

\subsection*{بخش ه – سنجش سه روش دفاعی}

\textbf{۱) فرستادن وقفه‌های دستگاه به یک هستهٔ دیگر.} وقفه‌های دستگاه‌ها (شبکه، \lr{GPU}، دیسک) دیگر هستهٔ مهاجم را متوقف نمی‌کنند؛ پس بزرگ‌ترین بخش سیگنال از بین می‌رود. اما کامل نیست: وقفهٔ زمان‌سنج و وقفه‌های بین‌هسته‌ای (در ادامه) باقی می‌مانند؛ کار نرم‌افزاری مرتبط با وقفه‌ها (مثلاً پردازش بسته‌ها) ممکن است هنوز روی هستهٔ مهاجم اجرا شود؛ و کانال‌های دیگری مثل کش و پهنای باند حافظه باقی‌اند. هزینه: یک هستهٔ پرمشغله برای همهٔ وقفه‌ها (افت کارایی و تأخیر) و نیاز به چسباندن پردازه‌ها به هسته. \textbf{میزان جلوگیری: متوسط (کاهش قابل توجه ولی نه حذف کامل).}

\textbf{۲) کم‌دقت و نوفه‌دار کردن ساعت.} حمله به ساعت دقیق نیاز ندارد: شمارنده فقط تعداد تکرار حلقه در یک بازهٔ چندمیلی‌ثانیه‌ای را می‌شمارد و وقفه‌ها (از چند میکروثانیه تا میلی‌ثانیه) در این مقیاس ظاهر می‌شوند. یعنی ساعتی با دقت چند ده یا صد میکروثانیه هم هنوز همین کار را می‌کند. مهاجم می‌تواند ساعت جایگزین هم بسازد (با نخ شمارنده و \lr{SharedArrayBuffer}). این روش فقط وقتی اثر دارد که ساعت چنان خراب شود که مرورگر عملاً بی‌فایده شود. \textbf{میزان جلوگیری: کم (ضعیف)}. این روش بیشتر در برابر حمله‌هایی کارساز است که به ساعت دقیق نیاز دارند.

\textbf{۳) تزریق وقفه‌های ساختگی و تصادفی.} وقفه‌های اضافه، الگوی واقعی را با نوفه می‌پوشانند و دقت دسته‌بند را کم می‌کنند. اما نوفهٔ مستقل از قربانی را می‌توان با میانگین‌گیری از ردهای تکراری و فیلتر کردن تا حدی حذف کرد. برای اینکه الگو واقعاً پنهان شود، نوفه باید هم‌اندازهٔ سیگنال و از نظر شکل و مدت شبیه وقفه‌های واقعی باشد. هزینه: هدر دادن \lr{CPU} و انرژی و افزایش تأخیر. \textbf{میزان جلوگیری: متوسط، با هزینهٔ زیاد.}

\begin{center}\small
\small
\begin{tabular}{@{}>{\raggedright\arraybackslash}p{3.4cm}>{\raggedright\arraybackslash}p{4.4cm}>{\raggedright\arraybackslash}p{4.2cm}>{\raggedright\arraybackslash}p{2.3cm}@{}}
\toprule
  {\textbf{روش}} &
  {\textbf{مهم‌ترین نقطهٔ قوت}} &
  {\textbf{مهم‌ترین نقطهٔ ضعف}} &
  {\textbf{میزان جلوگیری}} \\
\midrule
  {۱) انتقال وقفه‌ها} &
  {حذف وقفه‌های دستگاه‌ها} &
  {زمان‌سنج، بین‌هسته‌ای، کش} &
  {متوسط} \\
  {۲) ساعت کم‌دقت} &
  {ساده و ارزان} &
  {حمله به ساعت دقیق نیاز ندارد} &
  {کم} \\
  {۳) وقفهٔ ساختگی} &
  {پنهان‌سازی الگوی واقعی} &
  {میانگین‌گیری و هزینهٔ زیاد} &
  {متوسط} \\
\bottomrule
\end{tabular}
\end{center}

\textbf{چرا وقفهٔ زمان‌سنج و وقفه‌های بین‌هسته‌ای را نمی‌توان به هستهٔ دیگر منتقل کرد؟}
\begin{itemize}
  \item \textbf{وقفهٔ زمان‌سنج:} هر هسته زمان‌سنج محلی خودش را دارد (مثلاً \lr{LAPIC timer}) و سیستم‌عامل برای زمان‌بندی روی همان هسته به آن نیاز دارد. سیستم‌عامل پیش از اجرای پردازه زمان‌سنج را تنظیم می‌کند تا پس از پایان سهم زمانی کنترل را پس بگیرد. اگر این وقفه به هستهٔ دیگر برود، هستهٔ اجراکنندهٔ پردازه هرگز پیش‌دستی نمی‌شود و سهم زمانی و زمان‌بندی از بین می‌رود.
  \item \textbf{وقفه‌های بین‌هسته‌ای (\lr{IPI}):} ماهیتاً پیامی از یک هسته به هستهٔ هدف مشخص است تا حالت همان هسته عوض شود؛ مثل ابطال \lr{TLB} همان هسته، بیدار کردن یا پیش‌دستی روی آن، یا اجرای یک تابع روی آن. منتقل کردن آن به هسته‌ای دیگر معنایش را از بین می‌برد. در مقابل، وقفهٔ دستگاه از کنترل‌گر وقفه به هر هسته‌ای قابل مسیریابی است.
\end{itemize}


\section{کنترل‌گر مدیریت برد (\lr{BMC}) و راه‌اندازی}

\subsection*{بخش آ – مسیر راه‌اندازی میزبان و جایگاه \lr{BMC}}

برنامهٔ راه‌اندازی (میان‌افزار، در \lr{ROM/EPROM}) هنگام روشن شدن یا راه‌اندازی دوباره اجرا می‌شود، همه‌چیز را مقداردهی می‌کند (ثبات‌های \lr{CPU}، کنترل‌گرهای دستگاه)، هسته را بارگذاری می‌کند و هسته نخستین پردازه را اجرا می‌کند؛ مثال لینوکس: \lr{BIOS $\rightarrow$ MBR bootloader $\rightarrow$ GRUB $\rightarrow$ Kernel init $\rightarrow$ INIT $\rightarrow$ Run level}. مراحل \lr{UEFI} هم \lr{SEC}، \lr{PEI}، \lr{DXE} و \lr{BDS} است. \lr{BMC} رایانهٔ کوچکی (معمولاً \lr{SoC} مبتنی بر \lr{ARM}) روی برد اصلی است که \textbf{در مسیری جدا} قرار دارد.

\begin{center}
\begin{tikzpicture}[
  font=\small, >=Stealth,
  box/.style={draw=navy, fill=lane1, rounded corners=3pt, minimum height=1.05cm, minimum width=5.7cm, align=center, inner sep=3pt, line width=0.6pt},
  bmc/.style={draw=accent, fill=lane2, rounded corners=3pt, minimum height=1.05cm, minimum width=5.7cm, align=center, inner sep=3pt, line width=0.6pt},
  fin/.style={draw=navy, fill=white, rounded corners=3pt, minimum height=1.05cm, minimum width=5.7cm, align=center, inner sep=3pt, line width=1.2pt, dashed},
  main/.style={->, line width=0.7pt, color=navy},
  side/.style={->, line width=0.7pt, color=accent, dashed},
  lbl/.style={font=\scriptsize, align=center, fill=white, inner sep=1.5pt}
]
\node[box] (h0) at (4.4, 0)    {\RL{\textbf{روشن شدن برق}}\\[-1pt]\RL{\scriptsize بررسی سلامت برق (\lr{power-good})}};
\node[box] (h1) at (4.4,-1.65) {\RL{\textbf{میان‌افزار \lr{UEFI/BIOS}}}\\[-1pt]\RL{\scriptsize \lr{SEC}، \lr{PEI}: مقداردهی \lr{CPU}، حافظه و تراشه‌ها}};
\node[box] (h2) at (4.4,-3.30) {\RL{\textbf{\lr{DXE} و \lr{BDS}}}\\[-1pt]\RL{\scriptsize بارگذاری درایورها، انتخاب دستگاه بوت}};
\node[box] (h3) at (4.4,-4.95) {\RL{\textbf{بوت‌لودر (\lr{MBR/GRUB})}}\\[-1pt]\RL{\scriptsize پیدا کردن و بارگذاری هسته در حافظه}};
\node[box] (h4) at (4.4,-6.60) {\RL{\textbf{هسته (\lr{Kernel init})}}\\[-1pt]\RL{\scriptsize راه‌اندازی زیرسیستم‌ها و سوار کردن ریشه}};
\node[box] (h5) at (4.4,-8.25) {\RL{\textbf{نخستین پردازه (\lr{init/PID 1})}}\\[-1pt]\RL{\scriptsize اجرای سرویس‌ها و پردازه‌های کاربر}};
\node[fin] (h6) at (4.4,-9.90) {\RL{\textbf{محیط اجرای پردازه‌های کاربر}}\\[-1pt]\RL{\scriptsize (هدف مسئله: میزبان به اینجا نمی‌رسد)}};
\foreach \a/\b in {h0/h1,h1/h2,h2/h3,h3/h4,h4/h5,h5/h6} \draw[main] (\a.south) -- (\b.north);
\node[bmc] (b0) at (-4.4, 0)    {\RL{\textbf{برق آماده‌به‌کار و بوت داخلی}}\\[-1pt]\RL{\scriptsize \lr{BMC} پیش از میزبان و مستقل از آن بالا می‌آید}};
\node[bmc] (b1) at (-4.4,-1.65) {\RL{\textbf{رابط مدیریتی}}\\[-1pt]\RL{\scriptsize \lr{IPMI/Redfish/Web} روی شبکهٔ اختصاصی یا مشترک}};
\node[bmc] (b2) at (-4.4,-3.30) {\RL{\textbf{نمایش وضعیت و ثبت فرمان‌ها}}\\[-1pt]\RL{\scriptsize وضعیت میزبان، بازنشانی، حسگرها، کنسول}};
\foreach \a/\b in {b0/b1,b1/b2} \draw[main,color=accent] (\a.south) -- (\b.north);
\draw[side] (b2.east) -- ++(1.55,0) |- (h0.west);
\node[lbl] at (0,-1.65) {\RL{فرمان روشن،}\\\RL{خاموش، بازنشانی}\\\RL{(خط سخت‌افزاری)}};
\draw[side] (h4.west) -| (b2.south);
\node[lbl] at (-1.4,-6.60) {\RL{وضعیت، حسگرها،}\\\RL{کنسول سریال}};
\node[font=\bfseries\small, color=navy] at (4.4, 0.95) {\RL{مسیر اصلی (میزبان)}};
\node[font=\bfseries\small, color=accent] at (-4.4, 0.95) {\RL{مسیر مدیریتی جداگانه (\lr{BMC})}};
\end{tikzpicture}
\end{center}

\textbf{توضیح تصویر:} مسیر اصلی (راست) از برق، میان‌افزار، بوت‌لودر و هسته به نخستین پردازه می‌رسد. \lr{BMC} (چپ) رایانهٔ مستقل با برق آماده‌به‌کار، میان‌افزار و رابط شبکهٔ خودش است و فقط از طریق خطوط سخت‌افزاری و پیوندهای جانبی (روشن/بازنشانی، حسگرها، کنسول سریال) با مسیر میزبان ارتباط دارد. به همین دلیل در تمام مراحل و حتی وقتی میزبان به محیط اجرای پردازه‌های کاربر نمی‌رسد، فعال می‌ماند.

\subsection*{بخش ب – آیا رابط مدیریتی \lr{BMC} پردازه‌ای در سیستم‌عامل میزبان است؟}

\textbf{خیر.} دلیل‌ها:
\begin{itemize}
  \item \lr{BMC} یک رایانهٔ کوچک مستقل روی برد اصلی است (معمولاً \lr{ARM-based SoC}) با \lr{CPU}، حافظه و میان‌افزار (و معمولاً سیستم‌عامل) خودش. پردازه در مفهوم درس یعنی برنامه‌ای که هستهٔ میزبان آن را ایجاد و زمان‌بندی می‌کند و از حافظهٔ میزبان استفاده می‌کند؛ رابط \lr{BMC} هیچ‌کدام را ندارد.
  \item برق آماده‌به‌کار و درگاه شبکهٔ اختصاصی (یا مشترک از طریق \lr{NC-SI}) دارد؛ پس حتی وقتی میزبان خاموش است یا هنوز هسته‌ای بارگذاری نشده، کار می‌کند. چیزی که پیش از وجود سیستم‌عامل کار می‌کند نمی‌تواند پردازهٔ همان سیستم‌عامل باشد.
  \item بر اساس صورت مسئله، رابط مدیریتی با مسیر مدیریتی جداگانه از میزبان است و حتی وقتی سیستم‌عامل میزبان اجرا نشده پاسخ می‌دهد. بنابراین نه از هستهٔ میزبان زمان می‌گیرد و نه به آن وابسته است.
\end{itemize}

\subsection*{بخش ج – آیا پاسخ‌گو بودن \lr{BMC} در حالی که میزبان به محیط پردازه‌ها نرسیده تناقض ایجاد می‌کند؟}

\textbf{خیر، تناقضی نیست}، چون \lr{BMC} و میزبان دو سامانهٔ جدا هستند.
\begin{itemize}
  \item \textbf{آنچه می‌توان نتیجه گرفت:} برق آماده‌به‌کار برقرار است، \lr{BMC} بوت شده و مسیر مدیریتی (شبکه و رابط) سالم است، و \lr{BMC} می‌تواند وضعیت سخت‌افزاری برد را بخواند (برق، حسگرها) و فرمان بازنشانی را ثبت کند.
  \item \textbf{آنچه نمی‌توان نتیجه گرفت:} \lr{BMC} اطلاعات خود را از سیگنال‌ها و حسگرهای سخت‌افزاری می‌گیرد؛ مثلاً می‌بیند میزبان «روشن است». این به معنای سلامت یا اجرای میان‌افزار، بوت‌لودر، هسته یا هر پردازه‌ای نیست. نمی‌توان نتیجه گرفت که میزبان به‌درستی بوت شده یا از کار نیفتاده است؛ فقط می‌توان گفت میزبان به محیط اجرای پردازه‌های کاربر نرسیده است و علت آن (میان‌افزار، بوت‌لودر، هسته یا سخت‌افزار) را باید جای دیگر پیدا کرد.
\end{itemize}

\subsection*{بخش د – بارگذاری هسته پس از ثبت درخواست بازنشانی}

پس از روشن شدن دوباره، ابتدا میان‌افزار (\lr{UEFI/BIOS}) اجرا می‌شود و سخت‌افزار را مقداردهی می‌کند، سپس بوت‌لودر (\lr{GRUB}) \textbf{هسته} را از دیسک در حافظه بارگذاری می‌کند و کنترل را به آن می‌دهد. \textbf{ثبت درخواست بازنشانی به معنای اجرای سیستم‌عامل میزبان توسط \lr{BMC} نیست}، زیرا:
\begin{itemize}
  \item \lr{BMC} فقط یک «فرمان» ثبت و ارسال می‌کند (فعال‌کردن خط بازنشانی یا چرخهٔ برق؛ معادل زدن دکمهٔ \lr{reset}) و کد یا حافظهٔ میزبان را اجرا نمی‌کند.
  \item پس از دریافت سیگنال، مسئولیت اجرا با سخت‌افزار و میان‌افزار میزبان است که از ابتدا شروع می‌کنند و فرایند بوت را طی می‌کنند. کد \lr{BMC} روی پردازندهٔ جداگانه‌ای اجرا می‌شود و اصلاً به میزبان وارد نمی‌شود.
  \item هسته‌ای که بارگذاری می‌شود همان است که روی دیسک میزبان قرار دارد و \lr{BMC} در انتخاب یا بارگذاری آن نقشی ندارد (مگر فرمان‌هایی مثل تعیین دستگاه بوت در دفعهٔ بعد، که باز هم فقط یک تنظیم برای میان‌افزار است).
\end{itemize}

\subsection*{بخش ه – میزبان به درخواست‌های پردازهٔ کاربر پاسخ می‌دهد، ولی رابط \lr{BMC} در دسترس نیست}

\textbf{نمی‌توان نتیجه گرفت که سیستم‌عامل میزبان از کار افتاده است.} سیستم‌عامل میزبان به \lr{BMC} وابسته نیست؛ پاسخ‌گو بودن پردازه‌های کاربر نشان می‌دهد که هسته، زمان‌بندی و مسیر شبکهٔ میزبان سالم است. نبودن رابط \lr{BMC} به یکی از این علت‌ها می‌تواند باشد: \lr{BMC} در حال راه‌اندازی دوباره یا به‌روزرسانی میان‌افزار است، خود \lr{BMC} (میان‌افزار یا سخت‌افزارش) خراب شده، شبکهٔ مدیریتی یا درگاه اختصاصی (یا مسیر \lr{NC-SI}) قطع است، یا تنظیمات دسترسی عوض شده است. یعنی خرابی \lr{BMC} فقط توانایی مدیریت از راه دور را کم می‌کند و نشانهٔ خرابی میزبان نیست.

\subsection*{بخش و – مقایسهٔ دو حالت}

\begin{center}\small
\small
\begin{tabular}{@{}>{\raggedright\arraybackslash}p{4.4cm}>{\raggedright\arraybackslash}p{5.3cm}>{\raggedright\arraybackslash}p{5.3cm}@{}}
\toprule
  {\textbf{حالت}} &
  {\textbf{مسیری که کار می‌کند}} &
  {\textbf{مسیری که ممکن است از کار افتاده باشد}} \\
\midrule
  {\lr{BMC} در دسترس، سیستم‌عامل میزبان در دسترس نیست} &
  {مسیر برون‌باند \lr{BMC}: دیدن وضعیت برق و حسگرها، روشن/خاموش‌کردن و بازنشانی، کنسول سریال از راه دور، ثبت رویدادها} &
  {مسیر درون‌باند میزبان: هسته، پردازه‌های کاربر، سرویس‌هایی مثل \lr{SSH}، شبکهٔ میزبان؛ (سیستم‌عامل پایین است یا اصلاً بالا نیامده)} \\
\addlinespace
  {\lr{BMC} در دسترس نیست، سیستم‌عامل میزبان در دسترس است} &
  {مسیر درون‌باند میزبان: سیستم‌عامل، پردازه‌ها و شبکهٔ میزبان؛ (به‌جز حالتی که کارت شبکه را مشترک با \lr{BMC} استفاده کنند و آن مسیر هم قطع شود)} &
  {مسیر مدیریتی برون‌باند: رابط \lr{BMC}، کنترل از راه دور برق و بازنشانی، کنسول، حسگرها؛ (برق و میزبان مشکلی ندارند)} \\
\bottomrule
\end{tabular}
\end{center}

\section{سیگنال‌ها در لینوکس}

\subsection*{بخش آ – رفتارهای پیش‌فرض پس از ورود سیگنال}

برای هر سیگنال فقط پنج رفتار پیش‌فرض ممکن است (راهنمای \code{signal(7)}):
\begin{itemize}
  \item \code{Term}: پایان‌دادن به پردازه (مثل \code{SIGTERM}، \code{SIGINT}، \code{SIGKILL}، \code{SIGALRM}).
  \item \code{Core}: پایان‌دادن به پردازه و ساختن فایل \lr{core dump} از حافظهٔ پردازه (مثل \code{SIGSEGV}، \code{SIGABRT}، \code{SIGFPE}).
  \item \code{Ign}: نادیده‌گرفتن سیگنال (مثل \code{SIGCHLD}، \code{SIGURG}).
  \item \code{Stop}: متوقف‌کردن پردازه (مثل \code{SIGSTOP}، \code{SIGTSTP}).
  \item \code{Cont}: ادامه‌دادن پردازهٔ متوقف‌شده (\code{SIGCONT}).
\end{itemize}

دو مورد را توضیح می‌دهیم:
\begin{itemize}
  \item \textbf{\code{Term} (مثلاً \code{SIGTERM}، شمارهٔ ۱۵):} هسته پردازه را بدون ذخیرهٔ وضعیت متوقف می‌کند و منابعش را (حافظه، فایل‌های باز) آزاد می‌کند. والد با \code{wait} وضعیت پایان را می‌بیند. اگر برنامه گردانندهٔ سیگنال نصب کرده باشد، به‌جای رفتار پیش‌فرض آن اجرا می‌شود (برای مثال برای بستن مرتب فایل‌ها). \code{SIGKILL} و \code{SIGSTOP} را نمی‌توان گرفت، مسدود یا نادیده گرفت.
  \item \textbf{\code{Core} (مثلاً \code{SIGSEGV}، شمارهٔ ۱۱):} وقتی پردازه به آدرس نامعتبر دسترسی می‌دهد، سخت‌افزار تله ایجاد می‌کند، هسته آن را به سیگنال تبدیل می‌کند و پردازه پایان می‌یابد و تصویر حافظه و ثبات‌ها در فایل \lr{core dump} برای اشکال‌زدایی ذخیره می‌شود (اسلاید فصل ۲: «\lr{Failure of an application can generate core dump file}»).
\end{itemize}
\src{تکمیلی (\code{signal(7)})، اشاره به \lr{core dump} و تله در اسلایدها.}

\subsection*{بخش ب – مؤلفه‌های مرتبط با سیگنال در \code{/proc/PID/status}}

این فایل چند فیلد دربارهٔ سیگنال‌ها دارد که همگی \textbf{نقاب بیتی شانزده‌دهی} هستند: بیت شمارهٔ \lr{n-1} روشن یعنی سیگنال شمارهٔ \lr{n} در آن مجموعه است. سه مؤلفهٔ اصلی:
\begin{itemize}
  \item \code{SigBlk} (\lr{Blocked}): سیگنال‌های \textbf{مسدودشده} (نقاب سیگنال). تحویل سیگنال‌های این مجموعه به تعویق می‌افتد تا نقاب برداشته شود. مثال: \code{SigBlk: 0000000000010000} یعنی بیت ۱۶ (سیگنال ۱۷، یعنی \code{SIGCHLD}) مسدود است.
  \item \code{SigIgn} (\lr{Ignored}): سیگنال‌هایی که پردازه \textbf{نادیده می‌گیرد} (رفتار \code{SIG\_IGN}).
  \item \code{SigCgt} (\lr{Caught}): سیگنال‌هایی که پردازه برایشان \textbf{گردانندهٔ سفارشی} ثبت کرده است (مثل \code{sigHandler()} در شکل ۱ صورت تمرین).
\end{itemize}

فیلدهای مرتبط دیگر: \code{SigPnd} و \code{ShdPnd} (سیگنال‌های \textbf{در انتظار} برای نخ و برای کل پردازه) و \code{SigQ} (تعداد سیگنال‌های در صف نسبت به سقف مجاز، مثلاً \code{0/31538}).
\src{تکمیلی (\code{proc(5)}).}

\subsection*{بخش ج – سیگنال‌های تله‌های فراخوانی نامعتبر و خطای ممیز شناور}

\begin{itemize}
  \item \textbf{فراخوانی سیستمی نامعتبر:} \code{SIGSYS} (شمارهٔ ۳۱ در \lr{x86-64}، «\lr{Bad system call}»)، رفتار پیش‌فرض \code{Core}. نکته: در لینوکس یک شمارهٔ فراخوانی ناموجود معمولاً فقط خطای \code{ENOSYS} برمی‌گرداند و سیگنال نمی‌فرستد. \code{SIGSYS} وقتی ارسال می‌شود که یک فیلتر \lr{seccomp} فراخوانی را رد کند (\code{SECCOMP\_RET\_TRAP}) یا در سیستم‌های یونیکس دیگر مثل \lr{BSD} و \lr{Solaris}.
  \item \textbf{خطا در محاسبات (از جمله شناور):} \code{SIGFPE} (شمارهٔ ۸، «\lr{Floating-point exception}»)، رفتار پیش‌فرض \code{Core}. با وجود نامش، علاوه بر خطاهای ممیز شناور (تقسیم بر صفر، سرریز و ...) برای تقسیم صحیح بر صفر یا سرریز صحیح نیز ارسال می‌شود. سازوکار آن همان تله است که اسلایدها توضیح می‌دهند (تقسیم بر صفر یک وقفهٔ نرم‌افزاری ناشی از خطاست): \lr{CPU} استثنا ایجاد می‌کند، هسته آن را به \code{SIGFPE} تبدیل و به پردازه ارسال می‌کند. نوع خطا از \code{si\_code} (مثل \code{FPE\_INTDIV}، \code{FPE\_FLTDIV}) مشخص می‌شود.
\end{itemize}
\src{تکمیلی (\code{signal(7)}، \code{sigaction(2)})؛ مفهوم تله در اسلاید.}

\subsection*{بخش د – چرا امکان ماسک‌کردن وقفه لازم است؟}

\begin{itemize}
  \item \textbf{حفاظت از بخش‌های بحرانی و جلوگیری از شرایط رقابتی:} اگر گرداننده در میانهٔ به‌روزرسانی یک ساختمان داده (مثلاً صف زمان‌بندی) اجرا شود و همان ساختمان داده را تغییر دهد، داده ناسازگار می‌شود. با ماسک‌کردن، اجرای اتمیک تضمین می‌شود (برای نمونه هسته هنگام تغییر ساختمان‌های مشترک با گرداننده وقفه را غیرفعال می‌کند، و برنامه با \code{sigprocmask} سیگنال‌ها را مسدود می‌کند).
  \item \textbf{جلوگیری از تودرتوشدن و بازورود:} هنگام اجرای گرداننده، همان وقفه (یا وقفه‌های با اولویت پایین‌تر) نباید دوباره آن را قطع کند؛ وگرنه سرریز پشته یا ناسازگاری وضعیت رخ می‌دهد. سیستم‌عامل ذخیره و بازیابی وضعیت را هم باید بدون مزاحمت انجام دهد (اسلاید: سیستم‌عامل ثبات‌ها و شمارندهٔ برنامه را ذخیره می‌کند).
  \item \textbf{مرحلهٔ راه‌اندازی و تغییر حالت:} هنگام بوت، پیش از تنظیم جدول وقفه و گردانندها، یا هنگام سوئیچ زمینه، دریافت وقفه خطرناک است.
  \item \textbf{اولویت‌بندی و کنترل بار:} وقفه‌های مهم‌تر روی کم‌اهمیت‌ها مقدم می‌شوند و طوفان وقفه (بار زیاد) را می‌توان مهار کرد. وقفه ازدست‌نرفته است و تا برداشتن نقاب در حالت انتظار می‌ماند (در سیگنال‌های استاندارد چند ارسال پشت‌سرهم یکی حساب می‌شوند). وقفه‌های حیاتی (مثل خطای سخت‌افزاری شدید) غیرقابل‌ماسک (\lr{NMI}) هستند تا هرگز نادیده گرفته نشوند.
\end{itemize}
\src{تکمیلی؛ اشاره به ذخیرهٔ وضعیت هنگام وقفه در اسلاید فصل ۱.}

\end{document}
